%=======================================================================
%  PE5 -- Practica Experimental Unidad V
%  Proyecto Integrador de Ingenieria de Requisitos (ISR-401) -- UTEQ
%  Equipo AHMRV -- Paralelo 4to "A"
%  Villafuerte Rosero, Huilcapi Leon, Rizzo Velez,
%  Macias Herrera, Arboleda Yanza, Alcivar Velez
%
%  COMPILACION (tres pasadas obligatorias):
%     pdflatex PE5_U5_PFC_ALCIVAR_ARBOLEDA_HUILCAPI_MACIAS_RIZZO_VILLAFUERTE
%     bibtex   PE5_U5_PFC_ALCIVAR_ARBOLEDA_HUILCAPI_MACIAS_RIZZO_VILLAFUERTE
%     pdflatex PE5_U5_PFC_ALCIVAR_ARBOLEDA_HUILCAPI_MACIAS_RIZZO_VILLAFUERTE
%     pdflatex PE5_U5_PFC_ALCIVAR_ARBOLEDA_HUILCAPI_MACIAS_RIZZO_VILLAFUERTE
%
%  DEPENDENCIAS NO INCLUIDAS EN UNA INSTALACION MINIMA:
%     texlive-lang-spanish   (spanish.ldf)
%     texlive-publishers     (IEEEtran.bst)
%     texlive-science        (siunitx)
%     texlive-pictures       (pgfplots, tikz)
%=======================================================================
\documentclass[11pt,letterpaper]{article}

\usepackage[utf8]{inputenc}
\usepackage[T1]{fontenc}
% Carga condicional: si falta el modulo de espanol, el documento compila igual
\IfFileExists{spanish.ldf}{\usepackage[spanish,es-tabla]{babel}}{}
\usepackage[margin=2.5cm]{geometry}
\usepackage{graphicx}
\usepackage{booktabs}
\usepackage{longtable}
\usepackage{array}
\usepackage{tabularx}
\usepackage{multirow}
\usepackage{xcolor}
\usepackage{colortbl}
\usepackage{caption}
\usepackage{float}
\usepackage{enumitem}
\usepackage{fancyhdr}
\usepackage{siunitx}
\usepackage{pgfplots}
\pgfplotsset{compat=1.18}
\usepackage{tikz}
\usetikzlibrary{shapes.geometric, arrows.meta, positioning}
\usepackage[hidelinks]{hyperref}
\usepackage{xurl}
\setlength{\emergencystretch}{3em}
\graphicspath{{figuras/}}
\sisetup{output-decimal-marker={,}}

\definecolor{grisclaro}{RGB}{224,224,224}
\definecolor{grissuave}{RGB}{243,243,243}
\definecolor{verdesuave}{RGB}{226,239,218}
\definecolor{rojosuave}{RGB}{252,232,232}

\newcolumntype{L}[1]{>{\raggedright\arraybackslash}p{#1}}
\newcolumntype{C}[1]{>{\centering\arraybackslash}p{#1}}
% Columna centrada vertical y horizontalmente: evita que el filete del
% identificador sobresalga por debajo en filas de varias lineas
\newcolumntype{M}[1]{>{\centering\arraybackslash}m{#1}}
\newcommand{\id}[1]{\texttt{#1}}
\pagestyle{fancy}
\fancyhf{}
\fancyhead[L]{\footnotesize PE5 --- Proyecto Integrador de Ingeniería de Requisitos}
\fancyhead[R]{\footnotesize ISR-401 --- UTEQ}
\fancyfoot[C]{\footnotesize Página \thepage}
\renewcommand{\headrulewidth}{0.4pt}

\begin{document}

%=======================================================================
% CARATULA
%=======================================================================
\begin{titlepage}
\centering
\vspace*{0.8cm}

{\large\textbf{UNIVERSIDAD TÉCNICA ESTATAL DE QUEVEDO}}\\[0.3cm]
{\normalsize Facultad de Ciencias de la Computación}\\[0.2cm]
{\normalsize Carrera de Software (Rediseño)}\\[1.4cm]

{\large\textbf{PRÁCTICA EXPERIMENTAL --- UNIDAD V}}\\[0.5cm]
{\Large\textbf{Integración, Métricas y Defensa del Proyecto Integrador}}\\[0.4cm]
{\normalsize Auditoría de calidad ISO/IEC 25010:2023 · Trazabilidad extremo a extremo · Requisitos de componentes de inteligencia artificial}\\[1.4cm]

{\normalsize\textbf{Sistema del Proyecto Fin de Curso}}\\[0.3cm]
{\large SIMPA --- Sistema Inteligente de Mantenimiento de Palma Africana}\\[0.2cm]
{\footnotesize Documento base: ERS/SRS del SIMPA, versión 1.1 final}\\[1.2cm]

{\normalsize\textbf{Integrantes}}\\[0.3cm]
\begin{tabular}{L{5.8cm} L{3.6cm} L{3.4cm}}
\toprule
\textbf{Integrante} & \textbf{Rol} & \textbf{Usuario de GitHub}\\
\midrule
Villafuerte Rosero Allan Noe & Analista líder & \texttt{AlanNVR}\\
Huilcapi León Denisses Fabiola & Documentadora & \texttt{huilcapi}\\
Rizzo Vélez Edson Nagib & Verificador & \texttt{erizzov-boop}\\
Macías Herrera Josthyn Esteban & Modelador & \texttt{jmaciasherr4}\\
Arboleda Yanza Francisco Javier & Apoyo modelado & \texttt{farboleday-wq}\\
Alcívar Vélez Anderson Adonis & Apoyo repositorio & \texttt{AdonisAlcivar}\\
\bottomrule
\end{tabular}\\[1.2cm]

{\normalsize\textbf{Asignatura:} Ingeniería de Requerimientos (ISR-401) --- 4.º nivel}\\[0.2cm]
{\normalsize\textbf{Docente:} Ing. Gleiston Cicerón Guerrero Ulloa, PhD}\\[0.2cm]
{\normalsize\textbf{Periodo académico:} 2026--2027 PPA}\\[0.2cm]
{\normalsize\textbf{Fecha:} 21/08/2026}\\[1.2cm]

{\normalsize\textbf{Repositorio del trabajo}}\\[0.2cm]
{\small\url{https://github.com/erizzov-boop/ISR401-PFC-ERS-AHMRV}}

\end{titlepage}

%=======================================================================
% RESUMEN BILINGUE
%=======================================================================
\section*{Resumen}
\addcontentsline{toc}{section}{Resumen}

Este informe integra las cinco unidades del Proyecto Fin de Curso de Ingeniería de Requerimientos
sobre el SIMPA, sistema de gestión y diagnóstico para el cultivo de palma aceitera desarrollado con
una organización real del cantón El Empalme. El trabajo audita la Especificación de Requisitos de
Software final con seis métricas derivadas de ISO/IEC 25010:2023, calculadas sobre conteos
programáticos de las propias fuentes del documento y no sobre estimaciones. La primera medición
detectó cuatro métricas por debajo de su valor de referencia: la especificación de casos de uso, la
cobertura de actores, la verificabilidad y la corrección. Las tres primeras se cerraron especificando
ocho casos de uso, reescribiendo treinta y siete criterios de aceptación en formato Dado--Cuando--Entonces
y declarando la distinción entre clase de usuario y parte interesada. La cuarta se reporta sin
corregir, porque mide un hecho histórico sobre la eficacia de la inspección anterior. El trabajo
cierra la trazabilidad extremo a extremo en setenta y tres filas con ocho eslabones, incorpora los
diagramas de flujo de datos que faltaban y especifica dieciocho requisitos para dos componentes de
inteligencia artificial, con sus umbrales de rendimiento, equidad y explicabilidad, su clasificación
de riesgo conforme al Reglamento (UE) 2024/1689 y su plan de monitoreo en producción.

\vspace{0.3cm}
\noindent\textbf{Palabras clave:} ingeniería de requisitos, métricas de calidad, ISO/IEC 25010,
trazabilidad extremo a extremo, requisitos de inteligencia artificial, equidad algorítmica.

\vspace{0.8cm}
\section*{Abstract}
\addcontentsline{toc}{section}{Abstract}

This report integrates the five units of the Requirements Engineering final course project on SIMPA,
a management and diagnosis system for oil palm cultivation developed with a real organisation in the
El Empalme canton. The work audits the final Software Requirements Specification using six metrics
derived from ISO/IEC 25010:2023, computed from programmatic counts over the document sources rather
than from estimates. The first measurement found four metrics below their reference value: use case
specification, actor coverage, verifiability and correctness. The first three were closed by
specifying eight use cases, rewriting thirty-seven acceptance criteria in Given--When--Then format,
and stating explicitly the distinction between user class and stakeholder. The fourth is reported
uncorrected, since it measures a historical fact about the effectiveness of the previous inspection.
The work closes end-to-end traceability across seventy-three rows with eight links, incorporates the
missing data flow diagrams, and specifies eighteen requirements for two artificial intelligence
components, including performance, fairness and explainability thresholds, risk classification under
Regulation (EU) 2024/1689, and a production monitoring plan.

\vspace{0.3cm}
\noindent\textbf{Keywords:} requirements engineering, quality metrics, ISO/IEC 25010, end-to-end
traceability, artificial intelligence requirements, algorithmic fairness.

\newpage

%=======================================================================
% INDICES
%=======================================================================
\tableofcontents
\newpage
\listoftables
\addcontentsline{toc}{section}{Índice de tablas}
\listoffigures
\addcontentsline{toc}{section}{Índice de figuras}
\newpage

%=======================================================================
% SECCIONES
%=======================================================================
\section{Introducción}
\label{sec:intro}

\subsection{Problema y contexto}

La palma aceitera es un cultivo perenne cuya rentabilidad depende de decisiones que se toman en
ventanas de tiempo estrechas y sobre información que hoy no se registra. El proyecto se desarrolla
con una organización productora del cantón El Empalme, provincia del Guayas, que explota
aproximadamente cien hectáreas distribuidas en seis lotes y emplea a sesenta y dos personas.

El trabajo de elicitación de las unidades anteriores documentó tres problemas que el sistema debe
atender y que conviene enunciar antes que cualquier consideración técnica, porque son los que
justifican el proyecto:

\begin{enumerate}[leftmargin=*, nosep]
  \item \textbf{La detección fitosanitaria depende de una visita quincenal.} Entre visitas de la
        asesoría técnica, una afectación avanza sin que nadie la registre. El personal de campo la
        observa pero no dispone de un medio para clasificarla ni para dejar constancia.
  \item \textbf{La penalización por fruta verde se decide sobre un juicio visual.} La planta
        extractora aplica un descuento cuando el porcentaje de fruta verde supera un umbral, y la
        decisión de cortar se toma con una estimación hecha a ojo en el momento del corte.
  \item \textbf{La liquidación semanal se calcula a mano.} La evidencia \id{EV-13} documenta que el
        cálculo de la remuneración del personal consume una jornada administrativa completa cada
        semana y se realiza sobre anotaciones en papel.
\end{enumerate}

Ninguno de los tres es un problema de software. Son problemas de operación agrícola cuya solución
pasa por registrar información que hoy se pierde, y esa distinción es la que gobierna todo el
documento: el SIMPA no automatiza porque se pueda, sino donde la falta de registro tiene una
consecuencia económica medible sobre la organización o sobre las personas que trabajan en ella.

\subsection{El sistema}

El SIMPA ---Sistema Inteligente de Mantenimiento de Palma Africana--- es una aplicación móvil con
panel web de administración que cubre siete bloques funcionales: gestión de la estructura productiva,
planificación y seguimiento de labores, registro operativo de campo, polinización asistida,
diagnóstico asistido por inteligencia artificial, calidad de la fruta y relación con la extractora, y
reportes, estimación e histórico.

La versión final del ERS especifica 42 requisitos funcionales, 19 no funcionales, 10 restricciones de
diseño y 8 requisitos legales derivados de la Ley Orgánica de Protección de Datos Personales, más 18
requisitos específicos de los dos componentes de inteligencia artificial que esta unidad incorpora.

\subsection{Objetivos de esta práctica}

\noindent\textbf{Objetivo general.} Integrar, auditar y defender el proyecto: aplicar métricas de
calidad al ERS final, cerrar la trazabilidad extremo a extremo, especificar los requisitos de los
componentes de inteligencia artificial y sostener las decisiones de ingeniería ante el tribunal.

\noindent\textbf{Objetivos específicos.}

\begin{enumerate}[leftmargin=*, nosep]
  \item Construir un instrumento de auditoría con seis métricas calculadas sobre conteos verificables
        de las fuentes del propio documento.
  \item Corregir toda métrica por debajo de su valor de referencia y volver a medir, presentando el
        par antes/después.
  \item Cerrar la trazabilidad extremo a extremo, incorporando los eslabones de proceso, estado y
        criterio de aceptación que la cadena anterior no tenía.
  \item Especificar dos componentes de inteligencia artificial con sus requisitos de rendimiento,
        equidad y explicabilidad, su clasificación de riesgo y su plan de monitoreo.
  \item Consolidar la versión final del ERS y documentar el aporte individual verificable.
\end{enumerate}

\subsection{Estructura del documento}

La \S\ref{sec:metodologia} describe el proceso de ingeniería de requisitos seguido a lo largo de las
cinco unidades y justifica las decisiones metodológicas. La \S\ref{sec:ers} presenta la estructura
del ERS final y sus requisitos clave. La \S\ref{sec:uml} recorre los modelos con lectura
interpretativa. La \S\ref{sec:validacion} expone la inspección formal, sus defectos y la
re-inspección de esta unidad. La \S\ref{sec:gestion} trata la línea base, el control de cambios y la
matriz de trazabilidad final. La \S\ref{sec:ia} especifica los componentes de inteligencia
artificial. La \S\ref{sec:metricas} presenta la auditoría de calidad con la aritmética de cada
métrica. La \S\ref{sec:retro} recoge la retrospectiva y la \S\ref{sec:conclusiones} cierra con cinco
conclusiones, una por unidad.


% =====================================================================
\section{Metodología de ingeniería de requisitos}
\label{sec:metodologia}

\subsection{Marco normativo aplicado}

El ERS se estructura conforme a ISO/IEC/IEEE 29148:2018, que define el contenido de una
especificación de requisitos de software y las características exigibles a un requisito individual y
al conjunto \cite{iso29148}. La calidad del producto documental se interpreta con el modelo de
ISO/IEC 25010:2023, del que la \S\ref{sec:metricas} deriva seis métricas operativas
\cite{iso25010}. El encuadre de procesos de ciclo de vida proviene de ISO/IEC/IEEE 15288:2023 para
sistemas y de ISO/IEC/IEEE 12207:2017 para software \cite{iso15288, iso12207}, que sitúan la
definición de requisitos de las partes interesadas y el análisis de requisitos del sistema dentro de
los procesos técnicos.

Los fundamentos conceptuales siguen a Pohl \cite{pohl2025} y al SWEBOK v4.0 \cite{swebok4}, y la
lista de verificación de calidad de la especificación se apoya en el temario CPRE Foundation Level
del IREB \cite{ireb31}. La gestión del proyecto integrador se organiza con los principios del PMBOK
7.ª edición \cite{pmbok7}.

Para los componentes de inteligencia artificial concurren dos marcos adicionales: el Reglamento
(UE) 2024/1689, que clasifica los sistemas por nivel de riesgo e impone obligaciones proporcionales
\cite{aiact2024}, y la Ley Orgánica de Protección de Datos Personales del Ecuador, que rige el
tratamiento de los datos personales que el sistema recoge \cite{lopdp2021}.

\subsection{Fundamento teórico}

\subsubsection*{Validación frente a verificación}

La distinción es anterior a cualquier técnica y condiciona qué se puede concluir de una inspección.
Verificar es comprobar que el documento está bien construido conforme a un estándar; validar es
comprobar que especifica lo que las partes interesadas necesitan \cite{pohl2025}. Un ERS puede
superar la verificación y fallar la validación: sus requisitos serían completos, consistentes y
verificables, y describirían un sistema que nadie pidió.

ISO/IEC/IEEE 29148:2018 sitúa la validación en su cláusula 6.4 y enumera las características
exigibles a un requisito individual ---necesario, apropiado, no ambiguo, completo, singular,
factible, verificable, correcto y conforme--- y al conjunto ---completo, consistente, factible,
comprensible, capaz de ser validado \cite{iso29148}. Las seis métricas de la
\S\ref{sec:metricas} operacionalizan seis de esas características; las restantes no se midieron
porque no admiten un conteo objetivo sin introducir juicio del auditor.

\subsubsection*{Técnicas de validación y el método Fagan}

Las técnicas de validación se ordenan por grado de formalidad: revisión informal, recorrido
(\emph{walkthrough}), inspección formal, prototipado y listas de verificación \cite{pohl2025,
swebok4}. El método de Fagan es la inspección formal canónica y aporta tres elementos que las
técnicas menos formales no tienen: roles diferenciados con responsabilidad distinta, preparación
individual obligatoria antes de la reunión, y prohibición expresa de discutir soluciones durante la
sesión \cite{fagan1976}.

Los tres elementos responden al mismo problema. Sin roles, todos leen igual y detectan lo mismo. Sin
preparación individual, el primer comentario ancla al grupo. Y sin la prohibición de discutir
soluciones, la sesión deriva hacia el diseño y deja de recorrer el documento. Las métricas de
inspección ---densidad de defectos, distribución por tipo y severidad, tasa de detección por
inspector y esfuerzo--- permiten evaluar el proceso y no solo su resultado.

\subsubsection*{Gestión de requisitos: línea base, trazabilidad y cambio}

La gestión de la configuración de un ERS descansa en la noción de línea base: una versión aprobada y
congelada que sirve de referencia y a partir de la cual todo cambio se tramita \cite{swebok4}. Sin
línea base no hay control de cambios posible, porque no existe el estado respecto del cual algo
cambia.

La trazabilidad se distingue en \emph{pre-traceability} ---del requisito hacia su origen--- y
\emph{post-traceability} ---del requisito hacia los artefactos que lo realizan--- \cite{gotel1994}.
Ambas direcciones responden preguntas distintas: la primera, por qué existe un requisito; la segunda,
qué hay que tocar si cambia. La matriz de la \S\ref{sec:gestion} las implementa como columnas hacia
atrás y hacia adelante de una misma cadena.

El proceso de cambio ---solicitud formal, análisis de impacto, decisión de un órgano colegiado,
implementación y cierre--- sigue el modelo de control integrado de cambios del PMBOK
\cite{pmbok7}. La literatura de trazabilidad advierte de que el análisis de impacto estimado sin
recorrer la matriz es la causa más frecuente de subestimación del costo de un cambio
\cite{clelandhuang2012}.

\subsubsection*{Herramientas CASE y trazabilidad en entornos ágiles}

Las herramientas CASE de gestión de requisitos se clasifican por su posición en el ciclo de vida
---\emph{upper}, \emph{lower}, integradas--- y por el conjunto de funcionalidades esperables:
repositorio, atributos tipados, trazabilidad multinivel, \emph{baselining}, control de cambios,
colaboración e integración con el control de versiones \cite{wiegers2013}.

En entornos ágiles la trazabilidad no desaparece: se reconstruye por convención sobre el mensaje de
\emph{commit} y la vinculación de ramas a la incidencia \cite{clelandhuang2012}. El formato de
historia de usuario \cite{cohn2004} y el criterio de aceptación en notación Gherkin \cite{smart2014}
cumplen el papel que en la especificación documental cumplen la ficha de requisito y su criterio de
verificación. El proyecto mantiene ambas representaciones, y la \S\ref{sec:gestion} cuantifica el
costo de esa convivencia.

\subsubsection*{Ingeniería de requisitos para sistemas con aprendizaje automático}

Especificar un componente cuyo comportamiento se induce de datos plantea problemas que la ingeniería
de requisitos clásica no aborda. El resultado es probabilístico, de modo que un criterio de
aceptación binario no aplica y hay que sustituirlo por una métrica con umbral y conjunto de
evaluación declarado. La calidad depende del conjunto de datos tanto como del modelo, lo que convierte
las propiedades de los datos en requisitos. Y aparecen categorías de calidad nuevas ---explicabilidad
y ausencia de discriminación--- que ningún modelo de calidad tradicional contemplaba
\cite{vogelsang2024, chazette2021}.

A esos requisitos técnicos se suman los normativos. El Reglamento (UE) 2024/1689 adopta un enfoque
basado en riesgo: prohíbe determinadas prácticas, impone obligaciones reforzadas a los sistemas de
alto riesgo enumerados en su Anexo III, y establece obligaciones de transparencia cuando el sistema
interactúa directamente con personas \cite{aiact2024}. La normativa ecuatoriana de protección de
datos personales rige el tratamiento de los datos que alimentan el modelo y los derechos del titular
sobre ellos \cite{lopdp2021}. La \S\ref{sec:ia} aplica ambos marcos al SIMPA y declara la
clasificación resultante con su razonamiento.

\subsection{Proceso seguido a lo largo de las cinco unidades}

\begin{table}[H]
\centering\footnotesize
\caption{Proceso de ingeniería de requisitos por unidad}
\label{tab:proceso}
\begin{tabular}{|C{1.4cm}|L{4.0cm}|L{4.6cm}|L{4.6cm}|}
\hline
\rowcolor{grisclaro}\textbf{Unidad} & \textbf{Actividad} & \textbf{Técnicas empleadas} & \textbf{Artefacto producido}\\ \hline
I & Fundamentos y planificación & Análisis del dominio, identificación de partes interesadas, matriz poder/interés & Plan de elicitación y mapa de partes interesadas\\ \hline
II & Elicitación y análisis & Entrevista semiestructurada, cuestionario, observación de campo, análisis de documentos & Trece evidencias (\id{EV-01} a \id{EV-13}) y catálogo de requisitos brutos\\ \hline
III & Especificación y modelado & Fichas de requisito conforme a 29148, casos de uso, modelado i*, UML, diagramas de flujo de datos & ERS de las Entregas 1B y 2A; modelos UML e i*\\ \hline
IV & Validación y gestión & Inspección formal tipo Fagan, Change Control Board simulado, trazabilidad en Jira, línea base con Git & Registro de 33 defectos, tres RFC, acta del CCB, línea base \id{baseline-v1.1}\\ \hline
V & Integración y auditoría & Auditoría con métricas de 25010, re-inspección, especificación de componentes de IA & Instrumento de auditoría, matriz extremo a extremo, \S9 del ERS\\ \hline
\end{tabular}
\end{table}

\subsection{Decisiones metodológicas y su justificación}

Cuatro decisiones condicionaron el resultado y merecen justificarse, porque en cada una había una
alternativa razonable que se descartó.

\subsubsection*{Elicitación con instrumentos mixtos y no solo con entrevistas}

La elicitación combinó entrevista semiestructurada con la administración, cuestionario al personal de
campo y observación directa. La razón es una asimetría documentada en \id{EV-07}: la administración
declara con soltura los procesos formales, pero el personal de campo ---sesenta y dos personas con
competencia digital declarada como baja--- no responde bien al formato de entrevista. Basar la
elicitación solo en entrevistas habría producido un catálogo de requisitos que refleja cómo la
administración cree que funciona el campo, y no cómo funciona.

El cuestionario ampliado \id{EV-12}, con $n=62$, es el que aportó los dos datos que más restricciones
generaron: que el 11,3\,\% del personal no dispone de teléfono inteligente y que el 74,2\,\% no tiene
señal permanente en el lote. Ambos datos contradijeron suposiciones que el equipo había dado por
válidas y obligaron a especificar \id{RF-35} (registro delegado) y \id{RD-02} (operación sin
conexión).

\subsubsection*{Priorización con tres técnicas y no con una}

La priorización combina MoSCoW, el modelo de Kano y el cálculo WSJF. MoSCoW por sí solo produce
categorías sin orden interno: veinticuatro requisitos \emph{Must} no dicen cuál construir primero.
Kano aporta la distinción entre lo que satisface y lo que se da por supuesto. WSJF aporta el orden
dentro de la categoría, ponderando valor de negocio, urgencia, reducción de riesgo y tamaño.

El resultado de combinarlas no fue redundante. \id{RF-35} ---el registro delegado para quien no tiene
teléfono--- encabeza el orden de implementación con un WSJF de 7,00 pese a no ser el requisito de
mayor valor de negocio: lo que lo eleva es la reducción de riesgo, porque sin él el 11,3\,\% del
personal quedaría fuera del sistema y el registro de avance sería incompleto por diseño.

\subsubsection*{Inspección formal y no revisión informal}

La Unidad IV empleó el método Fagan con roles diferenciados, preparación individual previa registrada
con marca temporal y prohibición expresa de discutir soluciones durante la sesión \cite{fagan1976}.
La alternativa ---una revisión conjunta del documento--- es más barata y detecta menos: sin
preparación individual, la sesión se convierte en una lectura colectiva donde el primer comentario
ancla al resto del grupo.

El resultado empírico sostiene la decisión: treinta y tres defectos únicos sobre sesenta y tres
páginas efectivas, con una densidad de 0,52 defectos por página, y ningún rol por debajo del mínimo
de seis defectos preparados.

\subsubsection*{Formato BDD para los criterios de aceptación}

La Unidad V reescribió treinta y siete criterios de verificación en formato Dado--Cuando--Entonces
\cite{smart2014}. La alternativa era conservarlos en prosa, que es como estaban y que no es
incorrecto. La razón del cambio es que un criterio en prosa admite que dos evaluadores lleguen a
veredictos distintos sobre el mismo resultado, mientras que la plantilla obliga a fijar contexto,
evento y resultado observable. La verificabilidad, tal como la define ISO/IEC/IEEE 29148:2018, no es
que exista un criterio sino que el criterio permita decidir \cite{iso29148}.

Conviene precisar el alcance de ese cambio, porque es fácil sobreinterpretarlo: no se especificó
ningún requisito nuevo ni se modificó ningún umbral. Cambió la forma de enunciar la condición, no su
contenido.


% =====================================================================
\section{Especificación de requisitos final}
\label{sec:ers}

\subsection{Estructura del documento conforme a ISO/IEC/IEEE 29148:2018}

\begin{table}[H]
\centering\footnotesize
\caption{Correspondencia entre las secciones del ERS y ISO/IEC/IEEE 29148:2018}
\label{tab:estructura-ers}
\begin{tabular}{|C{1.2cm}|L{5.0cm}|L{4.2cm}|C{3.6cm}|}
\hline
\rowcolor{grisclaro}\textbf{§} & \textbf{Sección del ERS} & \textbf{Cláusula de 29148} & \textbf{Extensión}\\ \hline
1 & Introducción, propósito y alcance & 9.5.2 Introduction & 7 páginas\\ \hline
2 & Descripción general y partes interesadas & 9.5.3 Overall description & 12 páginas\\ \hline
3 & Requisitos específicos & 9.5.4 Specific requirements & 34 páginas\\ \hline
4 & Modelado UML e i* & 9.5.5 Verification (apoyo) & 11 páginas\\ \hline
5 & Priorización y trazabilidad & 6.5 Requirements traceability & 12 páginas\\ \hline
6 & Prototipo mínimo viable y limitaciones & 9.5.6 Supporting information & 6 páginas\\ \hline
7 & Conclusiones y trabajo pendiente & --- & 4 páginas\\ \hline
8 & Diagramas de flujo de datos & 9.5.5 (incorporada en la PE5) & 5 páginas\\ \hline
9 & Componentes de inteligencia artificial & 9.5.4 (incorporada en la PE5) & 8 páginas\\ \hline
\end{tabular}
\end{table}

\subsection{Catálogo final de requisitos}

\begin{table}[H]
\centering\footnotesize
\caption{Composición del catálogo de requisitos del ERS final}
\label{tab:catalogo}
\begin{tabular}{|L{5.6cm}|C{2.0cm}|C{2.4cm}|L{4.4cm}|}
\hline
\rowcolor{grisclaro}\textbf{Tipo} & \textbf{Cantidad} & \textbf{Rango} & \textbf{Observación}\\ \hline
Requisitos funcionales & 42 & \id{RF-01}--\id{RF-42} & 24 \emph{Must}, 16 \emph{Should}, 2 \emph{Could}\\ \hline
Requisitos no funcionales & 19 & \id{RNF-01}--\id{RNF-19} & Mapeados a las nueve características de ISO/IEC 25010:2023\\ \hline
Restricciones de diseño & 10 & \id{RD-01}--\id{RD-10} & Impuestas por el entorno o la contraparte\\ \hline
Requisitos legales & 8 & \id{RL-01}--\id{RL-08} & Derivados de la LOPDP\\ \hline
Requisitos de IA & 18 & \id{RF-IA-01}--\id{RNF-IA-12} & Incorporados en la Unidad V\\ \hline
\textbf{Total auditado} & \textbf{61} & --- & RF + RNF; las restricciones y los requisitos legales se auditan aparte\\ \hline
\end{tabular}
\end{table}

\subsection{Requisitos clave y por qué lo son}

De los cuarenta y dos requisitos funcionales, cinco concentran la propuesta de valor del sistema y
son los que conviene poder defender uno por uno.

\begin{table}[H]
\centering\footnotesize
\caption{Requisitos clave del SIMPA y su justificación en la evidencia}
\label{tab:req-clave}
\begin{tabular}{|C{1.4cm}|L{4.2cm}|L{4.4cm}|L{4.4cm}|}
\hline
\rowcolor{grisclaro}\textbf{ID} & \textbf{Requisito} & \textbf{Evidencia que lo sustenta} & \textbf{Por qué es clave}\\ \hline
\id{RF-07} & Diagnóstico de plaga o enfermedad por imagen & \id{EV-02}: la asesoría técnica visita cada quince días & Es la funcionalidad que reduce la ventana de detección de dos semanas a inmediata\\ \hline
\id{RF-22} & Estimación del porcentaje de fruta verde con alerta preventiva & \id{EV-04}: la extractora penaliza a partir del 3\,\% & Traslada la decisión de cortar desde el juicio visual a una medición anterior al despacho\\ \hline
\id{RF-28} & Registro de avance por unidad de labor & \id{EV-05}, \id{EV-06} & Es la entrada de la que dependen la liquidación, el indicador de desempeño y los reportes\\ \hline
\id{RF-35} & Registro delegado del avance & \id{EV-12}: el 11,3\,\% no tiene teléfono inteligente & Sin él, una parte del personal quedaría fuera del sistema por su equipamiento\\ \hline
\id{RF-37} & Cálculo de la remuneración semanal & \id{EV-13}: el cálculo manual consume una jornada semanal & Es el proceso de mayor frecuencia y mayor costo administrativo de la organización\\ \hline
\end{tabular}
\end{table}

\subsection{Catálogo completo de requisitos funcionales}

\begin{longtable}{|M{1.8cm}|L{9.6cm}|M{2.4cm}|}
\caption{Catálogo de requisitos funcionales del ERS final, agrupados por bloque}
\label{tab:rf-completo}\\
\hline
\rowcolor{grisclaro}\textbf{ID} & \textbf{Nombre del requisito} & \textbf{MoSCoW}\\ \hline
\endfirsthead
\hline
\rowcolor{grisclaro}\textbf{ID} & \textbf{Nombre del requisito} & \textbf{MoSCoW}\\ \hline
\endhead
\multicolumn{3}{|>{\columncolor{grissuave}}l|}{\textbf{Bloque I - Gestión de la estructura productiva}}\\ \hline
\id{RF-01} & Autenticación y control de acceso por rol & Must\\ \hline
\id{RF-02} & Gestión de plantaciones y lotes & Must\\ \hline
\id{RF-03} & Gestión de personal y equipos de trabajo & Must\\ \hline
\id{RF-10} & Gestión de variedades de palma & Must\\ \hline
\multicolumn{3}{|>{\columncolor{grissuave}}l|}{\textbf{Bloque II - Planificación y seguimiento de labores}}\\ \hline
\id{RF-26} & Planificación semanal de labores con presupuesto & Must\\ \hline
\id{RF-27} & Arrastre de labor incompleta a la semana siguiente & Should\\ \hline
\id{RF-34} & Registro de solicitud de insumos asociada al plan & Could\\ \hline
\id{RF-36} & Catálogo codificado de labores con tarifa por unidad & Must\\ \hline
\id{RF-37} & Liquidación semanal de presupuesto contra ejecución & Must\\ \hline
\id{RF-38} & Registro de labor no presupuestada & Should\\ \hline
\id{RF-39} & Aprobación de la liquidación semanal & Should\\ \hline
\id{RF-35} & Registro delegado para personal sin dispositivo propio & Must\\ \hline
\multicolumn{3}{|>{\columncolor{grissuave}}l|}{\textbf{Bloque III - Registro operativo de campo}}\\ \hline
\id{RF-04} & Registro de labores agrícolas & Must\\ \hline
\id{RF-28} & Registro de avance por unidad de labor & Must\\ \hline
\id{RF-29} & Consulta de avance acumulado y estimación de remuneración & Should\\ \hline
\id{RF-05} & Registro de monitoreo fitosanitario & Must\\ \hline
\id{RF-30} & Reporte de incidencia desde campo con evidencia fotográfica y georreferencia & Must\\ \hline
\id{RF-31} & Lista de verificación de pendientes por visita & Should\\ \hline
\id{RF-06} & Seguimiento de etapas del cultivo & Should\\ \hline
\id{RF-40} & Exportación de los datos personales de la persona titular & Must\\ \hline
\id{RF-41} & Rectificación de datos con conservación en bitácora & Must\\ \hline
\id{RF-42} & Supresión de datos personales con disociación del histórico productivo & Must\\ \hline
\id{RF-11} & Registro de análisis de suelo y foliar & Should\\ \hline
\id{RF-17} & Registro y consulta de datos climáticos & Should\\ \hline
\multicolumn{3}{|>{\columncolor{grissuave}}l|}{\textbf{Bloque IV - Polinización asistida}}\\ \hline
\id{RF-13} & Registro del proceso de polinización & Must\\ \hline
\id{RF-14} & Conteo georreferenciado de flores polinizadas & Must\\ \hline
\id{RF-15} & Registro y visualización de recorridos del personal & Should\\ \hline
\id{RF-16} & Evaluación de desempeño del personal & Should\\ \hline
\multicolumn{3}{|>{\columncolor{grissuave}}l|}{\textbf{Bloque V - Diagnóstico asistido por inteligencia artificial}}\\ \hline
\id{RF-07} & Detección de plagas y enfermedades por análisis de imagen & Must\\ \hline
\id{RF-08} & Diagnóstico de deficiencias nutricionales por análisis de imagen & Must\\ \hline
\id{RF-09} & Ficha técnica de plaga con recomendación de control & Should\\ \hline
\id{RF-21} & Clasificación de madurez del racimo por desprendimiento & Must\\ \hline
\id{RF-12} & Generación de alertas tempranas & Must\\ \hline
\multicolumn{3}{|>{\columncolor{grissuave}}l|}{\textbf{Bloque VI - Calidad de la fruta y relación con la extractora}}\\ \hline
\id{RF-22} & Estimación y alerta preventiva de porcentaje de fruta verde & Must\\ \hline
\id{RF-23} & Registro de la calificación de calidad emitida por la extractora & Should\\ \hline
\id{RF-24} & Trazabilidad de la calificación hasta la cuadrilla responsable & Should\\ \hline
\id{RF-25} & Control de la ventana entre corte y entrega & Should\\ \hline
\id{RF-33} & Registro de la guía de remisión del despacho & Could\\ \hline
\multicolumn{3}{|>{\columncolor{grissuave}}l|}{\textbf{Bloque VII - Reportes, estimación e histórico}}\\ \hline
\id{RF-18} & Estimación de producción y rendimiento & Must\\ \hline
\id{RF-32} & Compartición de la estimación de cosecha con la planta extractora & Should\\ \hline
\id{RF-19} & Generación de reportes & Must\\ \hline
\id{RF-20} & Historial de actividades, diagnósticos y eventos & Should\\ \hline
\end{longtable}

\noindent
La distribución resultante ---24 \emph{Must}, 16 \emph{Should} y 2 \emph{Could}--- concentra el
57\,\% del catálogo en la categoría de cumplimiento obligatorio. Es una proporción alta y conviene
justificarla en lugar de presentarla como dato: tres de los \emph{Must} (\id{RF-40} a \id{RF-42})
implementan obligaciones legales que no admiten priorización, y cuatro más (\id{RF-36} a \id{RF-39})
cubren el proceso de liquidación semanal, que la organización ejecuta sin excepción todas las
semanas. Un catálogo con menos \emph{Must} habría exigido declarar opcional alguna de esas dos cosas.

\subsection{Requisitos no funcionales por característica de calidad}

\begin{longtable}{|M{1.8cm}|L{3.6cm}|L{8.4cm}|}
\caption{Requisitos no funcionales del ERS final y su característica de calidad}
\label{tab:rnf-completo}\\
\hline
\rowcolor{grisclaro}\textbf{ID} & \textbf{Característica} & \textbf{Descripción abreviada}\\ \hline
\endfirsthead
\hline
\rowcolor{grisclaro}\textbf{ID} & \textbf{Característica} & \textbf{Descripción abreviada}\\ \hline
\endhead
\id{RNF-01} & Adecuación funcional & La clasificación de madurez del racimo alcanza una exactitud >= 85\% y una exhaustividad >= 80\% sobre la clase ``verde'' en el conjunto de prueb\\ \hline
\id{RNF-02} & Adecuación funcional & La detección de plagas alcanza una exactitud >= 80\% sobre las cinco afectaciones más frecuentes del cultivo.\\ \hline
\id{RNF-03} & Eficiencia de desempeño & Las operaciones de registro y consulta responden en <= 3 s para el percentil 95, con hasta 20 sesiones concurrentes. El valor se deriva de la pob\\ \hline
\id{RNF-04} & Eficiencia de desempeño & El diagnóstico por imagen entrega resultado en <= 10 s con ancho de banda >= 5 Mbps.\\ \hline
\id{RNF-05} & Eficiencia de desempeño & El registro GPS captura la ubicación con precisión <= 10 m y período de muestreo <= 30 s.\\ \hline
\id{RNF-06} & Compatibilidad & La exportación hacia la planta extractora se emite en formato CSV o JSON conforme al esquema documentado, sin campos de identificación personal.\\ \hline
\id{RNF-07} & Compatibilidad & La aplicación móvil opera en Android 9.0 o superior; el panel web, en las dos últimas versiones estables de Chrome, Edge y Firefox.\\ \hline
\id{RNF-08} & Interacción de usuario & Una persona sin experiencia previa completa el registro de labor y el análisis de imagen con una tasa de éxito >= 90\% tras una inducción <= 10 m\\ \hline
\id{RNF-09} & Interacción de usuario & La totalidad de la interfaz se presenta en español, con etiquetas correspondientes a la terminología del glosario y sin abreviaturas técnicas no \\ \hline
\id{RNF-10} & Fiabilidad & El servicio de respaldo presenta una disponibilidad mensual >= 99\%, equivalente a una indisponibilidad <= 7 h 12 min (720 h 1\%).\\ \hline
\id{RNF-11} & Fiabilidad & La aplicación opera sin conexión y sincroniza el 100\% de los registros pendientes al restablecerse la conectividad, sin pérdida ni duplicación.\\ \hline
\id{RNF-12} & Fiabilidad & Ante pérdida de señal GPS, el sistema conserva el registro parcial y lo marca como de precisión degradada, sin descartar la jornada.\\ \hline
\id{RNF-13} & Seguridad & Las credenciales se almacenan mediante función de derivación de clave con sal; la totalidad del tráfico emplea TLS 1.2 o superior.\\ \hline
\id{RNF-14} & Seguridad & Los datos de geolocalización se cifran en reposo y su acceso queda registrado en bitácora con identificación de la persona usuaria, fecha y motiv\\ \hline
\id{RNF-15} & Mantenibilidad & El sistema presenta arquitectura modular con cobertura de pruebas automatizadas >= 60\% en la capa de lógica de negocio.\\ \hline
\id{RNF-16} & Interacción de usuario & Explicabilidad - véase la especificación ampliada en la §3.3.1.\\ \hline
\id{RNF-17} & Flexibilidad (reemplazabilidad) & La migración del servicio de inferencia a un proveedor alternativo no requiere modificar el cliente móvil, al mediar una interfaz de servicio des\\ \hline
\id{RNF-18} & Flexibilidad (adaptabilidad) & Los umbrales de madurez, de porcentaje de fruta verde y de ventana corte-entrega son parametrizables por variedad y por lote sin recompilar la ap\\ \hline
\id{RNF-19} & Seguridad física (safety) & Toda recomendación de control fitosanitario emitida por el sistema (RF-09) incluye, de forma no omisible, el equipo de protección personal exigid\\ \hline
\end{longtable}

\subsection{Restricciones de diseño y requisitos legales}

\begin{table}[H]
\centering\footnotesize
\caption{Restricciones de diseño y su origen}
\label{tab:rd}
\begin{tabular}{|C{1.6cm}|L{6.8cm}|L{6.2cm}|}
\hline
\rowcolor{grisclaro}\textbf{ID} & \textbf{Restricción} & \textbf{Origen}\\ \hline
\id{RD-01} & Cliente de campo sobre Android 9.0 o superior & El 88,7\,\% del personal dispone de teléfono inteligente propio (\id{EV-12})\\ \hline
\id{RD-02} & Operación sin conexión con sincronización diferida & El 74,2\,\% no tiene señal permanente en el lote (\id{EV-12})\\ \hline
\id{RD-05} & La inferencia se ejecuta dentro de la frontera del sistema & Restricción de conectividad y de costo de servicio externo\\ \hline
\id{RD-06} & Interfaz adaptada a competencia digital baja & Brecha declarada entre administración y personal de campo (\id{EV-07})\\ \hline
\id{RD-10} & Registro delegado para quien no dispone de dispositivo & El 11,3\,\% del personal (\id{EV-12})\\ \hline
\end{tabular}
\end{table}

\noindent
Se presentan cinco de las diez restricciones, las que condicionan decisiones de arquitectura; el
catálogo completo reside en la Tabla 52 del ERS. Las cinco comparten una propiedad que conviene
destacar: ninguna proviene de una preferencia técnica del equipo. Todas se derivan de un dato del
contexto que la elicitación midió.

\begin{table}[H]
\centering\footnotesize
\caption{Requisitos legales derivados de la LOPDP y su implementación}
\label{tab:rl}
\begin{tabular}{|C{1.4cm}|C{1.6cm}|L{5.6cm}|L{5.6cm}|}
\hline
\rowcolor{grisclaro}\textbf{ID} & \textbf{Artículo} & \textbf{Obligación} & \textbf{Requisito que la implementa}\\ \hline
\id{RL-04} & Art. 12 & Derecho de acceso de la persona titular a sus datos & \id{RF-40}, exportación por la propia titular\\ \hline
\id{RL-05} & Art. 13 & Derecho de rectificación & \id{RF-41}, corrección con conservación en bitácora\\ \hline
\id{RL-06} & Art. 14 & Derecho de eliminación & \id{RF-42}, supresión con disociación del histórico\\ \hline
\end{tabular}
\end{table}

\noindent
Los tres requisitos de la columna derecha no existían en la Entrega 2A. El ERS enunciaba las
obligaciones en su mapeo legal sin que ninguna función del sistema las satisficiera, y la inspección
de la Unidad IV lo registró como el defecto \id{DEF-12}. Su reparación exigía especificar tres
requisitos funcionales nuevos, de modo que no admitía corrección editorial y se tramitó ante el
Change Control Board como \id{RFC-03}.

\subsection{Requisitos no funcionales y modelo de calidad}

Los diecinueve requisitos no funcionales se organizan según las nueve características de calidad de
producto de ISO/IEC 25010:2023 \cite{iso25010}. La cobertura de las nueve se declaró en la Entrega 2A
y resultó no sostenerse: la inspección de la Unidad IV detectó que se empleaba ``Portabilidad'' como
característica de primer nivel ---cuando en la edición 2023 sus subcaracterísticas dependen de
Flexibilidad--- y que ninguna característica de Seguridad física tenía requisito asociado. Ambos
hallazgos se tramitaron ante el Change Control Board como \id{RFC-02} y se cerraron con la
reclasificación de \id{RNF-17} y la especificación de \id{RNF-19}.

El caso es ilustrativo de una diferencia que conviene retener: declarar conformidad con una norma y
aplicar su estructura correctamente son dos cosas distintas, y solo la segunda es verificable.

\section{Modelos UML y de procesos}
\label{sec:uml}

\subsection{Cobertura del modelado}

\begin{table}[H]
\centering\footnotesize
\caption{Modelos elaborados y su función en la especificación}
\label{tab:modelos}
\begin{tabular}{|L{4.2cm}|C{1.6cm}|L{8.8cm}|}
\hline
\rowcolor{grisclaro}\textbf{Modelo} & \textbf{Cantidad} & \textbf{Qué responde que los demás no responden}\\ \hline
Diagrama de contexto & 1 & Dónde está la frontera del sistema y qué queda fuera\\ \hline
Modelado i* (SD y SR) & 2 & Qué depende cada actor de los demás, y por qué querría usar el sistema\\ \hline
Casos de uso & 18 & Qué secuencia de pasos produce cada resultado, con sus alternativas y excepciones\\ \hline
Diagrama de clases refinado & 1 & Qué entidades del dominio existen y cómo se relacionan\\ \hline
Diagramas de secuencia & 3 & Cómo colaboran los objetos en los flujos más críticos\\ \hline
Diagrama de actividad & 1 & Cómo se encadena el ciclo semanal completo\\ \hline
Diagramas de estados & 2 & Qué ciclo de vida siguen el racimo y la alerta\\ \hline
Diagramas de flujo de datos & 2 & Qué proceso transforma qué dato y dónde se almacena\\ \hline
Componentes y despliegue & 2 & Cómo se reparte la funcionalidad y dónde se ejecuta\\ \hline
\end{tabular}
\end{table}

\subsection{Lectura interpretativa de los modelos principales}

\subsubsection*{Casos de uso: el hallazgo de la auditoría}

El diagrama general identifica dieciocho casos de uso, pero hasta esta unidad solo diez tenían
especificación textual. Los ocho restantes ---\id{CU-11} a \id{CU-18}--- existían como nodo del
diagrama y como referencia en la matriz, sin flujo, sin precondición y sin excepciones.

El vacío no lo detectó ninguna revisión de contenido: lo detectó la métrica M1b al dividir casos
especificados entre casos identificados y obtener 55,6\,\%. Es un ejemplo de por qué una auditoría
cuantitativa aporta algo que la lectura no aporta: un lector que recorre el documento de principio a
fin ve diez casos de uso bien especificados y no echa en falta los que no están.

Los ocho se especificaron en esta unidad con el mismo formato que los diez anteriores. Dos merecen
comentario. \id{CU-16} (generar alerta temprana) tiene como actor primario al propio sistema, que
actúa sin intervención humana; su flujo alternativo B evita la duplicación de alertas equivalentes,
que es el defecto que produciría un panel inservible por saturación. \id{CU-18} (explicar una
decisión automática) es incluido obligatoriamente por los tres casos de uso con decisión automática,
de modo que la explicabilidad no es una función que el usuario invoque sino una condición de la
propia emisión del resultado.

\subsubsection*{Diagramas de estados: dos entidades y por qué solo dos}

Se modelan el racimo y la alerta. La elección no es arbitraria ni exhaustiva: son las dos entidades
del dominio cuyo ciclo de vida determina una consecuencia económica y cuya transición no es trivial.

El estado \emph{Degradado} del racimo materializa el hallazgo de \id{EV-04} sobre la ventana entre
corte y entrega: el material híbrido tolera un intervalo mayor por su contenido de antioxidantes,
mientras que el \emph{guineensis} se deshidrata en dos o tres días. Sin ese estado, la ventana
corte--entrega de \id{RF-25} sería una regla sin objeto sobre el que aplicarse.

El estado \emph{Descartada} de la alerta corresponde a los falsos positivos del modelo de diagnóstico
y alimenta el ciclo de reentrenamiento previsto en \id{RNF-15} y \id{RNF-17}. Es el punto donde el
modelo de estados y los componentes de inteligencia artificial se tocan: un falso positivo no es solo
un error que se descarta, es un dato de entrenamiento.

\subsubsection*{Diagramas de flujo de datos: por qué se incorporaron ahora}

Los diagramas de flujo de datos se incorporaron en esta unidad, no en la de modelado. La razón es
concreta: la cadena de trazabilidad extremo a extremo exige enlazar cada requisito con el proceso que
lo realiza, y sin diagrama de flujo de datos esa columna quedaría vacía.

La decisión de diseño que conviene declarar es la correspondencia uno a uno entre los siete procesos
del nivel 1 y los siete bloques funcionales de la \S3.2 del ERS. No es una coincidencia del dominio:
se adoptó deliberadamente para que la introducción del modelo no obligara a reclasificar ningún
requisito ya especificado y para que un cambio en un requisito se propague por una sola ruta. La
alternativa ---descomponer el sistema por criterio propio del análisis de flujo--- habría producido
un modelo posiblemente más fiel pero con dos taxonomías paralelas del mismo sistema, que es
exactamente el mecanismo que generó los defectos de consistencia detectados en la Unidad IV.

\subsubsection*{Modelado i*: lo que aportó y lo que no}

Los diagramas SD y SR modelan las dependencias entre actores antes de que exista un sistema. Su
aporte concreto al proyecto fue detectar que la planta extractora es una parte interesada con poder
alto e interés medio, cuya calificación de calidad es una dependencia externa que el sistema no
controla pero de la que depende una función completa (\id{RF-23}, \id{RF-24}).

Lo que no aportó, y conviene decirlo: el modelado i* no produjo requisitos que la elicitación no
hubiera producido ya. Su valor fue de encuadre y de verificación cruzada, no de descubrimiento.


% =====================================================================
\section{Validación de requisitos}
\label{sec:validacion}

\subsection{Inspección formal de la Unidad IV}

\begin{table}[H]
\centering\footnotesize
\caption{Parámetros y resultados de la inspección formal tipo Fagan}
\label{tab:inspeccion}
\begin{tabular}{|L{6.4cm}|L{8.6cm}|}
\hline
\rowcolor{grisclaro}\textbf{Parámetro} & \textbf{Valor}\\ \hline
Documento inspeccionado & ERS/SRS del SIMPA v1.0 (Entrega 2A)\\ \hline
Alcance & \S1 a \S7; apéndices A--E excluidos\\ \hline
Páginas efectivas & 63 (pp. 8--70)\\ \hline
Roles & Moderador, Lector, Inspector 1 e Inspector 2, con acumulación declarada de dos roles en un integrante\\ \hline
Conformación del equipo en la inspección & Tres integrantes (conformación vigente en la Unidad IV)\\ \hline
Defectos únicos consolidados & 33\\ \hline
Distribución por severidad & 4 críticos, 24 mayores, 5 menores\\ \hline
Distribución por tipo & Consistencia 11, completitud 6, verificabilidad 5, ambigüedad 4, trazabilidad 4, factibilidad 3\\ \hline
Densidad de defectos & 0,52 defectos por página\\ \hline
Esfuerzo & 14,5 horas-persona; 0,44 horas-persona por defecto detectado\\ \hline
\end{tabular}
\end{table}

\noindent
Una precisión sobre la conformación del equipo antes de leer los resultados. La inspección se ejecutó
con los tres integrantes que componían el equipo de práctica experimental en la Unidad IV, lo que
obligó a acumular dos de los cuatro roles exigidos en un mismo participante; la acumulación se
declaró en su momento y no se disimula aquí. Con posterioridad a la inspección, el docente reagrupó
los equipos de práctica experimental conforme a los equipos del Proyecto Fin de Curso, con lo que
este equipo pasó a seis integrantes y la acumulación dejó de ser necesaria para trabajos ulteriores.
Los parámetros de la tabla anterior, la tasa de detección por rol y las métricas de inspección
corresponden a la ejecución original y no se recalculan: miden un hecho ocurrido en una fecha
determinada y con unos participantes determinados. Recalcularlos sobre la conformación actual
supondría atribuir a seis personas una detección que realizaron tres.

\vspace{0.3cm}
\noindent
Dos lecturas de esos números importan más que los números mismos.

La primera es la distribución por tipo. Once de los treinta y tres defectos son de consistencia y
ninguno es de sintaxis o de redacción. El documento sabía qué debía especificar; lo que fallaba era
que las mismas cifras se declaraban de forma distinta en secciones distintas. Es la firma de un
documento crecido por adición, sin una pasada de reconciliación final, y no la de un documento mal
escrito.

La segunda es el esfuerzo por defecto. Detectar un defecto costó menos de media hora-persona, cifra
que sostiene el argumento de Fagan de que la inspección es más barata que la corrección tardía
\cite{fagan1976}. El defecto \id{DEF-02} ---la contradicción entre 21 y 19 requisitos \emph{Must}---
habría llegado al diseño y habría producido dos planificaciones incompatibles.

\subsection{Corrección y gestión del cambio}

De los veintiocho defectos de severidad crítica y mayor, veinticinco se corrigieron directamente
sobre las fuentes y tres se tramitaron ante el Change Control Board por exceder el alcance de una
corrección editorial: su reparación consistía en añadir alcance al sistema, no en enmendar un texto.

\begin{table}[H]
\centering\footnotesize
\caption{Defectos derivados al Change Control Board}
\label{tab:def-ccb}
\begin{tabular}{|C{1.6cm}|C{1.6cm}|L{5.4cm}|L{6.0cm}|}
\hline
\rowcolor{grisclaro}\textbf{Defecto} & \textbf{RFC} & \textbf{Naturaleza y contenido} & \textbf{Decisión}\\ \hline
\id{DEF-13} & \id{RFC-01} & Alcance: \id{RF-36} y \id{RF-37} son \emph{Must} y carecían de historia de usuario & Aprobado por unanimidad\\ \hline
\id{DEF-16} & \id{RFC-02} & Calidad: la cobertura de las nueve características de ISO/IEC 25010:2023 no se sostenía & Aprobado por unanimidad\\ \hline
\id{DEF-12} & \id{RFC-03} & Normativa: \id{RL-04} a \id{RL-06} enunciaban obligaciones de la LOPDP sin requisito que las implementara & Aprobado con condiciones\\ \hline
\end{tabular}
\end{table}

\noindent
La deliberación sobre \id{RFC-03} merece registrarse porque el Change Control Board aprobó un cambio
que \emph{empeora} una métrica del propio proyecto. Incorporar \id{RF-40} a \id{RF-42} eleva los
requisitos \emph{Must} de veintiuno a veinticuatro y reduce la cobertura declarada del prototipo del
38,1\,\% al 33,3\,\%. El argumento que decidió la votación fue que el cumplimiento de una obligación
legal no se negocia contra un indicador de avance, y que el riesgo de omitirla recae sobre la
organización. Un órgano de control que solo aprobara lo que mejora sus indicadores no estaría
gobernando alcance sino administrando apariencia.

\subsection{Registro completo de defectos}

\begin{longtable}{|M{1.5cm}|L{2.4cm}|L{7.4cm}|M{1.9cm}|M{1.5cm}|}
\caption{Registro consolidado de los 33 defectos detectados en la inspección formal}
\label{tab:defectos}\\
\hline
\rowcolor{grisclaro}\textbf{ID} & \textbf{Sección} & \textbf{Descripción abreviada} & \textbf{Tipo} & \textbf{Sev.}\\ \hline
\endfirsthead
\hline
\rowcolor{grisclaro}\textbf{ID} & \textbf{Sección} & \textbf{Descripción abreviada} & \textbf{Tipo} & \textbf{Sev.}\\ \hline
\endhead
\id{DEF-01} & \S 3 (intro) & El encabezado de la \S 3 anuncia "9 restricciones de diseño"; la Tabla 64 y la \S 7.1 declaran 10 (RD-01 a RD-10). & Consistencia & Mayor\\ \hline
\id{DEF-02} & \S 5.1 vs \S 5.3.1 & La \S 5.1 declara una distribución MoSCoW de 21 Must, 16 Should y 2 Could (39 RF); la \S 5.3.1 declara 19 Must, 14 Should y 2 Could (35 RF). Las dos cifras no... & Consistencia & Crítico\\ \hline
\id{DEF-03} & \S 5.3.3, Anexo C & Ambos remiten a un archivo con "la tabla completa para los 35 requisitos funcionales", cuando el ERS especifica 39. & Consistencia & Mayor\\ \hline
\id{DEF-04} & \S 5.4, Tabla 69, \S 7.1, Anexo D & El número de filas de la matriz de trazabilidad se declara alternativamente como 52 y como 48 dentro del mismo documento; el diccionario define TR-01 a TR-48. & Consistencia & Mayor\\ \hline
\id{DEF-05} & \S 2.2 vs \S 3.2 & La \S 2.2 agrupa las funciones del producto en seis bloques; la \S 3.2 desarrolla siete (Bloque I a Bloque VII). & Consistencia & Mayor\\ \hline
\id{DEF-06} & Tabla 52 (RD-01) & RD-01 se justifica con "el personal ya dispone de teléfono inteligente propio (EV-05, EV-06, EV-07)". La \S 2.7 declara esa suposición invalidada por EV-12... & Consistencia & Crítico\\ \hline
\id{DEF-07} & Tabla 52 (RD-02) & RD-02 se justifica con "ninguna persona respondiente declaró señal permanente (EV-10)"; la \S 2.6 y la Tabla 7 reportan que el 25,8\% sí declara señal siempre... & Consistencia & Mayor\\ \hline
\id{DEF-08} & Tabla 78 vs Tabla 24 & El catálogo de evidencias atribuye a EV-13 el origen de RF-29, pero la ficha de RF-29 declara como origen EV-06 y EV-07 únicamente. & Consistencia & Mayor\\ \hline
\id{DEF-09} & \S 5.2, \S 5.4.1 (Tabla 70) & Se afirma que los 19 requisitos Must tienen historia de usuario y caso de uso (100\%). Si los Must son 21, la cobertura declarada es falsa y quedan requisitos... & Consistencia & Crítico\\ \hline
\id{DEF-10} & Tabla 68 & La tabla dice presentar los diez requisitos ordenados de forma descendente por WSJF, pero omite RF-05 y RF-30, a los que la Tabla 71 asigna WSJF 4,80, superior... & Consistencia & Mayor\\ \hline
\id{DEF-11} & \S 5.3.3 & La lectura del resultado afirma que RF-35 y RF-10 encabezan el orden con WSJF 7,00, pero RF-10 no figura en la Tabla 68; su valor solo aparece en la Tabla 71. & Consistencia & Mayor\\ \hline
\id{DEF-12} & Tabla 51 (\S 3.4) & RL-04, RL-05 y RL-06 enuncian obligaciones de la LOPDP (exportación de datos personales, rectificación con bitácora, supresión al término de la relación... & Completitud & Mayor\\ \hline
\id{DEF-13} & \S 5.2 & RF-36 y RF-37, clasificados Must en sus fichas, no tienen historia de usuario ni criterio de aceptación entre HU-01 y HU-19, pese a que la \S 5.2 declara cubrir... & Completitud & Crítico\\ \hline
\id{DEF-14} & \S 1.2.3 & La lista de funcionalidades principales omite el ciclo de liquidación semanal (RF-36 a RF-39), que la \S 7.2 califica como proceso central de la organización. & Completitud & Menor\\ \hline
\id{DEF-15} & Tabla 73 (L-06) & L-06 declara que las sesiones de walkthrough son inferiores a las tres exigidas y que "se declara el número efectivo"; ese número no se declara en ninguna... & Completitud & Mayor\\ \hline
\id{DEF-16} & \S 3.3 & Se afirma que los RNF cubren las nueve características de ISO/IEC 25010:2023, pero se emplea "Portabilidad" como característica de primer nivel (en la edición... & Completitud & Mayor\\ \hline
\id{DEF-17} & Tabla 39 (RF-12) & El requisito incluye "condiciones anómalas" como disparador de alerta sin definir qué constituye una condición anómala; el criterio de verificación solo cubre... & Verificabilidad & Mayor\\ \hline
\id{DEF-18} & Tabla 34 (RF-16) & No se define la fórmula, las variables ni la escala del "indicador de desempeño". El criterio de verificación ("el sistema calcula un indicador") es... & Verificabilidad & Mayor\\ \hline
\id{DEF-19} & Tabla 45 (RF-18) & La estimación de producción carece de tolerancia, margen de error o valor de referencia, por lo que ningún resultado puede declararse incorrecto. & Verificabilidad & Mayor\\ \hline
\id{DEF-20} & Tabla 49 (RNF-10) & Una disponibilidad mensual $\geq$99\% equivale a un máximo de 7 h 12 min de indisponibilidad (720 h $\times$ 1\%), no a las 7 h 18 min declaradas. & Verificabilidad & Menor\\ \hline
\id{DEF-21} & Tabla 49 (RNF-03) & Se fijan 50 sesiones concurrentes sin justificación; la población declarada es de 62 personas con conectividad mayoritariamente intermitente (74,2\% sin señal... & Verificabilidad & Menor\\ \hline
\id{DEF-22} & Tabla 35 (RF-07) & La precondición exige que "la imagen cumple los requisitos mínimos de calidad", pero esos requisitos no se especifican en ninguna sección del documento. & Ambigüedad & Mayor\\ \hline
\id{DEF-23} & Tabla 40 (RF-22) & La descripción exige alertar cuando el porcentaje "se aproxime al umbral", sin cuantificar la proximidad; el criterio de verificación usa "alcance o supere",... & Ambigüedad & Mayor\\ \hline
\id{DEF-24} & Tabla 43 (RF-25) / Tabla 52 (RD-08) & "El valor recomendado para la variedad" no está tabulado en ninguna parte. La única cifra aparece en la leyenda de la Figura 10 ("24 h híbrido / menor en... & Ambigüedad & Mayor\\ \hline
\id{DEF-25} & Tabla 37 (RF-09) & Se exige recomendar la acción de control "según la etapa del insecto", pero las etapas del ciclo no están enumeradas ni acotadas; el criterio de verificación... & Ambigüedad & Menor\\ \hline
\id{DEF-26} & Tabla 69 vs Tabla 57 & La matriz enlaza RF-16, RF-17 y RF-20 a CU-05, pero la especificación textual de CU-05 declara únicamente RF-19 y RF-26. La trazabilidad no cierra en sentido... & Trazabilidad & Mayor\\ \hline
\id{DEF-27} & Tabla 69 vs Tabla 9 & La matriz asocia RF-29 y RF-36 a RF-39 con el prototipo MU-02, mientras que la Tabla 9 declara que MU-02 soporta únicamente RF-12, RF-19 y RF-26. & Trazabilidad & Mayor\\ \hline
\id{DEF-28} & \S 1.4 y Referencias & Se declara que el documento cita 31 fuentes primarias; las entradas [27], [28], [30] y [31] figuran en la lista de referencias sin aparecer citadas en el texto. & Trazabilidad & Mayor\\ \hline
\id{DEF-29} & Tabla 78 vs portada & EV-13 se fecha entre el 03 y el 07/08/2026, es decir, con recolección posterior a la fecha de versión del propio documento (03/08/2026). & Trazabilidad & Mayor\\ \hline
\id{DEF-30} & Tabla 49 (RNF-01, RNF-02) & Los valores objetivo se verifican sobre un conjunto etiquetado de $\geq$200 imágenes que L-05 declara inexistente, y el documento no incorpora en el alcance de la... & Factibilidad & Mayor\\ \hline
\id{DEF-31} & \S 3.2.4 (RF-14) / Tabla 49 (RNF-05) & RF-14 exige derivar el conteo de flores planta por planta a partir de marcaciones GPS y RNF-05 admite una precisión de hasta 10 m. El ERS no declara la... & Completitud & Mayor\\ \hline
\id{DEF-32} & Tabla 73 & Los identificadores de limitación no siguen orden correlativo (L-01...L-07, L-10, L-11, L-08, L-09) y L-07 se declara "Resuelta" pero permanece listada como... & Factibilidad & Menor\\ \hline
\id{DEF-33} & \S 6.1.1 (Tabla 71) & La tabla se titula "cobertura sobre los requisitos funcionales Must" pero incluye RF-20 y RF-15, que sus fichas clasifican como Should. El denominador de 19 y... & Factibilidad & Mayor\\ \hline
\end{longtable}

\subsection{Métricas de la inspección e interpretación}

\begin{table}[H]
\centering\footnotesize
\caption{Métricas de la inspección formal y su lectura}
\label{tab:met-inspeccion}
\begin{tabular}{|L{3.8cm}|L{3.6cm}|L{7.6cm}|}
\hline
\rowcolor{grisclaro}\textbf{Métrica} & \textbf{Valor} & \textbf{Qué revela del ERS}\\ \hline
Densidad de defectos & 0,52 def./página & Con más de un defecto cada dos páginas, el documento no falla en un punto aislado sino de forma distribuida: apunta a un problema de integración entre secciones redactadas por separado\\ \hline
Distribución por tipo & Consistencia 11; completitud 6; verificabilidad 5; ambigüedad 4; trazabilidad 4; factibilidad 3 & La consistencia concentra un tercio: el documento sabe qué especificar, pero las mismas cifras se declaran distinto en secciones distintas\\ \hline
Distribución por severidad & 4 críticos; 24 mayores; 5 menores & El predominio de la severidad mayor sobre la menor indica que los defectos son de contenido especificado y no de forma\\ \hline
Tasa de detección por rol & Inspector 2: 11; Inspector 1: 8; Moderador: 7; Lector: 7 & Los cuatro roles superan el mínimo de seis y ninguno concentra más del 33\,\%: el reparto por secciones funcionó\\ \hline
Esfuerzo & 14,5 h-persona; 0,44 h-persona por defecto & Detectar un defecto costó menos de media hora-persona, cifra que justifica repetir la inspección antes de la siguiente entrega\\ \hline
\end{tabular}
\end{table}

\subsection{Re-inspección de la Unidad V}

La métrica de corrección de la auditoría requiere medir cuántos defectos sobrevivieron a la
inspección. Esa medición exige una re-inspección sobre el documento ya corregido, y su resultado es
el siguiente.

\begin{longtable}{|C{1.6cm}|L{4.0cm}|L{5.4cm}|C{1.5cm}|L{2.8cm}|}
\hline
\rowcolor{grisclaro}\textbf{ID} & \textbf{Ubicación} & \textbf{Defecto residual} & \textbf{Sev.} & \textbf{Estado}\\ \hline
\endhead
\id{RES-01} & Apéndice, diccionario de identificadores & Declara el rango de requisitos funcionales como \id{RF-01} a \id{RF-39}, cuando \id{RFC-03} lo elevó a \id{RF-42} & Mayor & Corregido\\ \hline
\id{RES-02} & Apéndice, diccionario de identificadores & Declara el rango de no funcionales como \id{RNF-01} a \id{RNF-18}, cuando \id{RFC-02} incorporó \id{RNF-19} & Mayor & Corregido\\ \hline
\id{RES-03} & Tabla de recuentos de la \S5.1 & Declara 39 requisitos funcionales y 18 no funcionales, cifras anteriores a las RFC & Mayor & Corregido\\ \hline
\id{RES-04} & Apéndice C & \id{RNF-19} no figura en la tabla de priorización: es el único requisito del catálogo sin prioridad MoSCoW & Mayor & Corregido\\ \hline
\id{RES-05} & \S4, casos de uso & Ocho de los dieciocho casos de uso identificados carecen de especificación textual & Mayor & Corregido\\ \hline
\id{RES-06} & \S5.2, criterios de aceptación & Solo 24 de 61 requisitos tienen criterio en formato Dado--Cuando--Entonces & Mayor & Corregido\\ \hline
\id{RES-07} & \S2.3 y fichas & Dos partes interesadas se contaban como clases de usuario sin ser actor de ningún requisito & Menor & Corregido\\ \hline
\caption{Defectos residuales detectados en la re-inspección de la Unidad V}
\label{tab:residuales}
\end{longtable}

\subsection{Por qué escaparon: análisis de causa}

Los siete defectos residuales no se distribuyen al azar, y su patrón dice algo sobre el
procedimiento de inspección más que sobre el documento.

\textbf{Cuatro se concentran en los apéndices y en los conteos globales} (\id{RES-01} a \id{RES-04}).
El alcance declarado de la inspección de la Unidad IV fue \S1 a \S7, con los apéndices excluidos de
forma expresa. La exclusión era razonable ---los apéndices son material de referencia--- pero dejó
fuera precisamente el lugar donde se acumulan los recuentos que las tres RFC modificaron. La causa no
es descuido: es que el alcance se fijó antes de saber qué iban a cambiar las RFC, y nadie lo revisó
después.

\textbf{Dos son de completitud del conjunto, no de corrección de sus elementos} (\id{RES-05},
\id{RES-06}). La lista de verificación empleada preguntaba si cada elemento revisado era correcto, no
si el conjunto estaba completo. Un inspector que recorre diez casos de uso bien especificados
responde ``sí cumple'' a todas las preguntas, y ninguna le obliga a contar cuántos deberían existir.

\textbf{Uno es un defecto de la propia auditoría} (\id{RES-07}). La primera medición de la cobertura
de actores contó ocho actores mezclando la tabla de clases de usuario con el mapa de partes
interesadas, y arrojó 75\,\% cuando el valor real era 100\,\%. Se corrigió el conteo y se añadió al
ERS una nota que declara la distinción, de modo que la métrica no vuelva a medirse mal. Es el único
de los siete que no era un defecto del ERS sino del instrumento que lo mide, y se documenta como tal
en lugar de retirarlo del registro.

\subsection{Acción correctiva de proceso}

Las siete correcciones están aplicadas, pero la conclusión operativa no es haberlas aplicado. Es que
la lista de verificación de la inspección debe modificarse en dos puntos para la siguiente entrega:
incluir los apéndices y los conteos globales dentro del alcance cuando el documento haya sufrido
cambios aprobados, e incorporar preguntas sobre la completitud del conjunto y no solo sobre la
corrección de sus elementos. Corregir siete defectos mejora este documento; cambiar la lista de
verificación mejora todos los siguientes.

\section{Gestión de requisitos y trazabilidad}
\label{sec:gestion}

\subsection{Línea base y control de versiones}

El ERS se versiona en el repositorio del equipo junto con sus fuentes \LaTeX{}, de modo que el
documento entregado se regenera por clonación y compilación. La línea base aprobada por el Change
Control Board al cierre de la Unidad IV se publicó como etiqueta anotada \id{baseline-v1.1}, y la
Unidad V cierra con la etiqueta \id{baseline-final}.

\begin{table}[H]
\centering\footnotesize
\caption{Historial de versiones del ERS}
\label{tab:versiones}
\begin{tabular}{|C{2.0cm}|C{2.2cm}|C{2.4cm}|L{7.4cm}|}
\hline
\rowcolor{grisclaro}\textbf{Versión} & \textbf{Rev. interna} & \textbf{Etiqueta Git} & \textbf{Contenido}\\ \hline
v1.0 & 3.0 & --- & Entrega 2A. Documento sometido a inspección formal\\ \hline
v1.1 & 3.1 & \id{baseline-v1.1}, \id{baseline-final} & Correcciones de la inspección, cambios aprobados por el CCB, y cierre de la Unidad V: auditoría, DFD, casos de uso, criterios BDD y componentes de IA\\ \hline
\end{tabular}
\end{table}

\noindent
La numeración merece una aclaración porque conviven dos. El ERS registra internamente sus revisiones
desde la Entrega 1A y llegó a la Entrega 2A como revisión interna 3.0; la asignatura designa a ese
mismo documento como v1.0. Se adopta la numeración de la asignatura como versión oficial, declarada
de forma idéntica en la portada del ERS, en su historial, en el \id{CHANGELOG.md} y en la etiqueta de
Git, y el historial del documento rotula su columna como \emph{revisión interna} para que la
coexistencia no produzca ambigüedad.

\subsection{Control de cambios}

Las tres solicitudes tramitadas ante el Change Control Board se recogen en la
\S\ref{sec:validacion}. El procedimiento seguido ---formulario con análisis de impacto derivado de la
matriz, deliberación con roles diferenciados, decisión motivada y acta firmada--- responde al modelo
de control integrado de cambios del PMBOK \cite{pmbok7}.

El análisis de impacto de cada solicitud se obtuvo recorriendo la matriz de trazabilidad y filtrando
por el requisito afectado, no por estimación. Esa diferencia es la que separa un análisis de impacto
de una conjetura informada: \id{RFC-03} declaró afectar a las filas \id{TR-03} y \id{TR-34} y añadir
\id{TR-53} a \id{TR-55}, y esos identificadores son verificables en el artefacto.

% =====================================================================
%  Sección de trazabilidad extremo a extremo — informe PE5
%  Requiere: longtable, array, xcolor, colortbl, float, y la macro \id
% =====================================================================

\subsection{Trazabilidad extremo a extremo}

\subsubsection*{Cadena de trazabilidad adoptada}

La trazabilidad de las entregas anteriores enlazaba la parte interesada con el requisito, el caso de
uso y el prototipo. Esa cadena bastaba para responder de dónde viene un requisito y qué pantalla lo
materializa, pero no para responder qué proceso lo ejecuta, sobre qué estado del dominio actúa ni
con qué condición se acepta. La Unidad~V extiende la cadena a ocho eslabones:

\begin{center}
\footnotesize
Fuente $\rightarrow$ Requisito $\rightarrow$ Caso de uso $\rightarrow$ Clase/Módulo $\rightarrow$
Proceso DFD $\rightarrow$ Estado $\rightarrow$ Criterio BDD $\rightarrow$ Caso de prueba
\end{center}

\noindent
Tres eslabones son nuevos y ninguno existía como columna antes de esta práctica:

\begin{itemize}[leftmargin=*, nosep]
  \item \textbf{Proceso DFD.} No existía porque el ERS no tenía diagramas de flujo de datos. Se
        incorporaron en la \S8 del ERS con siete procesos, uno por bloque funcional, de modo que la
        columna se deriva del bloque en que el propio documento ya ubicaba cada requisito y no
        introduce una clasificación nueva.
  \item \textbf{Estado.} Enlaza el requisito con la transición de la máquina de estados sobre la que
        actúa. El ERS modela el ciclo de vida de dos entidades: el racimo y la alerta.
  \item \textbf{Criterio BDD.} Enlaza el requisito con su condición de aceptación en formato
        Dado--Cuando--Entonces: \id{CA-01} a \id{CA-24} cuando procede de una historia de usuario y
        \id{CB-01} a \id{CB-37} cuando se deriva de la ficha del requisito.
\end{itemize}

\subsubsection*{Criterio de completitud de una traza}

Una fila se declara \textbf{completa} cuando tiene los seis eslabones exigibles: fuente, caso de uso,
módulo, proceso, criterio de aceptación y caso de prueba. Se declara \textbf{parcial} si falta uno y
\textbf{huérfana} si faltan dos o más.

El eslabón de estado no entra en ese conteo, y la razón conviene declararla porque afecta a la
lectura de la métrica. Solo dos entidades del dominio tienen ciclo de vida modelado; los requisitos
cuyo objeto no es un racimo ni una alerta ---gestión de personal, planificación de labores,
generación de reportes--- no tienen estado asociado. Exigirles uno obligaría a inventar máquinas de
estados sin sustento en el dominio para que una métrica saliera mejor. La columna declara
\emph{No aplica} en esas 34 filas, y esa declaración es información, no un vacío.

\subsubsection*{Matriz completa}

{\scriptsize
\begin{longtable}{|M{0.7cm}|L{1.9cm}|M{1.3cm}|M{1.0cm}|M{1.0cm}|M{1.1cm}|L{3.2cm}|M{1.1cm}|M{1.0cm}|M{1.3cm}|}
\caption{Matriz de trazabilidad extremo a extremo (73 filas). El asterisco indica que el proceso se deriva del caso de uso y no del bloque funcional, por tratarse de un requisito no funcional o de una restricción. En la columna de estado, \textbf{R} designa la máquina de estados del racimo y \textbf{A} la de la alerta.}
\label{tab:matriz-e2e}\\
\hline
\rowcolor{grisclaro}\textbf{N.º} & \textbf{Fuente} & \textbf{Requisito} & \textbf{CU} & \textbf{Módulo} & \textbf{Proceso} & \textbf{Estado} & \textbf{BDD} & \textbf{Prueba} & \textbf{Traza}\\ \hline
\endfirsthead
\hline
\rowcolor{grisclaro}\textbf{N.º} & \textbf{Fuente} & \textbf{Requisito} & \textbf{CU} & \textbf{Módulo} & \textbf{Proceso} & \textbf{Estado} & \textbf{BDD} & \textbf{Prueba} & \textbf{Traza}\\ \hline
\endhead
1 & Administrador & \id{RF-01} & CU-11 & M-01 & P-01 & No aplica & CA-01 & CP-01 & Completa\\ \hline
2 & Administrador & \id{RF-02} & CU-12 & M-01 & P-01 & No aplica & CA-02 & CP-02 & Completa\\ \hline
3 & Administrador & \id{RF-03} & CU-13 & M-01 & P-01 & No aplica & CA-03 & CP-03 & Completa\\ \hline
4 & Trabajador & \id{RF-04} & CU-04 & M-03 & P-03 & No aplica & CA-04 & CP-04 & Completa\\ \hline
5 & Trabajador & \id{RF-04} & CU-04 & M-03 & P-03 & No aplica & CA-04 & CP-04 & Completa\\ \hline
6 & Técnico fitosan. & \id{RF-05} & CU-15 & M-03 & P-03 & A: Abierta $\to$ En verificación & CA-05 & CP-05 & Completa\\ \hline
7 & Administrador & \id{RF-06} & CU-12 & M-03 & P-03 & R: En desarrollo & CB-01 & CP-06 & Completa\\ \hline
8 & Asistente admin. & \id{RF-07} & CU-01 & M-05 & P-05 & A: Abierta $\to$ Notificada & CA-06 & CP-07 & Completa\\ \hline
9 & Asistente admin. & \id{RF-07} & CU-01 & M-05 & P-05 & A: Abierta $\to$ Notificada & CA-06 & CP-07 & Completa\\ \hline
10 & Asesor técnico & \id{RF-08} & CU-02 & M-05 & P-05 & A: Abierta $\to$ Notificada & CA-07 & CP-08 & Completa\\ \hline
11 & Asesor técnico & \id{RF-08} & CU-02 & M-05 & P-05 & A: Abierta $\to$ Notificada & CA-07 & CP-08 & Completa\\ \hline
12 & Técnico fitosan. & \id{RF-09} & CU-01 & M-05 & P-05 & A: En verificación $\to$ Atendida & CB-02 & CP-09 & Completa\\ \hline
13 & Asesor técnico & \id{RF-10} & CU-12 & M-01 & P-01 & No aplica & CA-08 & CP-10 & Completa\\ \hline
14 & Asesor técnico & \id{RF-11} & CU-02 & M-03 & P-03 & No aplica & CB-03 & CP-11 & Completa\\ \hline
15 & Administrador & \id{RF-12} & CU-16 & M-05 & P-05 & A: (inicial) $\to$ Abierta & CA-09 & CP-12 & Completa\\ \hline
16 & Trabajador & \id{RF-12} & CU-16 & M-05 & P-05 & A: (inicial) $\to$ Abierta & CA-09 & CP-12 & Completa\\ \hline
17 & Jefe polinización & \id{RF-13} & CU-03 & M-04 & P-04 & R: Inflorescencia $\to$ Polinizado & CA-10 & CP-13 & Completa\\ \hline
18 & Jefe polinización & \id{RF-14} & CU-03 & M-04 & P-04 & R: Inflorescencia $\to$ Polinizado & CA-11 & CP-14 & Completa\\ \hline
19 & Administrador & \id{RF-15} & CU-03 & M-04 & P-04 & No aplica & CB-04 & CP-15 & Completa\\ \hline
20 & Administrador & \id{RF-16} & CU-05 & M-04 & P-04 & No aplica & CB-05 & CP-16 & Completa\\ \hline
21 & Administrador & \id{RF-17} & CU-05 & M-03 & P-03 & No aplica & CB-06 & CP-17 & Completa\\ \hline
22 & Administrador & \id{RF-18} & CU-17 & M-07 & P-07 & R: En desarrollo & CA-12 & CP-18 & Completa\\ \hline
23 & Administrador & \id{RF-19} & CU-05 & M-07 & P-07 & No aplica & CA-13 & CP-19 & Completa\\ \hline
24 & Administrador & \id{RF-20} & CU-05 & M-07 & P-07 & No aplica & CB-07 & CP-20 & Completa\\ \hline
25 & Téc. extractora & \id{RF-21} & CU-06 & M-05 & P-05 & R: En desarrollo $\to$ Maduro & CA-14 & CP-21 & Completa\\ \hline
26 & Téc. extractora & \id{RF-21} & CU-06 & M-05 & P-05 & R: En desarrollo $\to$ Maduro & CA-14 & CP-21 & Completa\\ \hline
27 & Administrador & \id{RF-22} & CU-06 & M-06 & P-06 & R: Maduro $\to$ Cosechado & CA-15 & CP-22 & Completa\\ \hline
28 & Administrador & \id{RF-23} & CU-10 & M-06 & P-06 & R: Despachado $\to$ Calificado & CB-08 & CP-23 & Completa\\ \hline
29 & Téc. extractora & \id{RF-24} & CU-10 & M-06 & P-06 & R: Calificado & CB-09 & CP-24 & Completa\\ \hline
30 & Téc. extractora & \id{RF-25} & CU-10 & M-06 & P-06 & R: Cosechado $\to$ Degradado & CB-10 & CP-25 & Completa\\ \hline
31 & Administrador & \id{RF-26} & CU-07 & M-02 & P-02 & No aplica & CA-16 & CP-26 & Completa\\ \hline
32 & Administrador & \id{RF-27} & CU-07 & M-02 & P-02 & No aplica & CB-11 & CP-27 & Completa\\ \hline
33 & Trabajador & \id{RF-28} & CU-04 & M-03 & P-03 & No aplica & CA-17 & CP-28 & Completa\\ \hline
34 & Trabajador & \id{RF-29} & CU-08 & M-03 & P-03 & No aplica & CB-12 & CP-29 & Completa\\ \hline
35 & Asistente admin. & \id{RF-30} & CU-09 & M-03 & P-03 & A: Abierta $\to$ En verificación & CA-18 & CP-30 & Completa\\ \hline
36 & Asistente admin. & \id{RF-31} & CU-09 & M-03 & P-03 & A: En verificación $\to$ Atendida & CB-13 & CP-31 & Completa\\ \hline
37 & Téc. extractora & \id{RF-32} & CU-17 & M-07 & P-07 & No aplica & CB-14 & CP-32 & Completa\\ \hline
38 & Téc. extractora & \id{RF-33} & CU-10 & M-06 & P-06 & R: Cosechado $\to$ Despachado & CB-15 & CP-33 & Completa\\ \hline
39 & Administrador & \id{RF-34} & CU-07 & M-02 & P-02 & No aplica & CB-16 & CP-34 & Completa\\ \hline
40 & Jefe de campo & \id{RF-35} & CU-04 & M-02 & P-02 & No aplica & CA-19 & CP-35 & Completa\\ \hline
41 & Téc. extractora & \id{RNF-01} & CU-06 & M-05 & P-05$^{*}$ & No aplica & CB-19 & CP-36 & Completa\\ \hline
42 & Trabajador & \id{RNF-11} & CU-04 & M-03 & P-03$^{*}$ & No aplica & CB-29 & CP-37 & Completa\\ \hline
43 & Jefe polinización & \id{RNF-05} & CU-03 & M-04 & P-04$^{*}$ & No aplica & CB-23 & CP-38 & Completa\\ \hline
44 & Trabajador & \id{RNF-08} & CU-04 & M-03 & P-03$^{*}$ & No aplica & CB-26 & CP-39 & Completa\\ \hline
45 & Trabajador & \id{RNF-14} & CU-03 & M-04 & P-04$^{*}$ & No aplica & CB-32 & CP-40 & Completa\\ \hline
46 & Asistente admin. & \id{RNF-16} & CU-18 & M-05 & P-05$^{*}$ & No aplica & CB-34 & CP-41 & Completa\\ \hline
47 & Trabajador & \id{RD-10} & CU-04 & M-03 & P-03$^{*}$ & No aplica & No aplica (restricción) & CP-42 & Completa\\ \hline
48 & Téc. extractora & \id{RD-07} & CU-06 & M-05 & P-05$^{*}$ & No aplica & No aplica (restricción) & CP-43 & Completa\\ \hline
49 & Administrador & \id{RF-36} & CU-08 & M-02 & P-02 & No aplica & CA-20 & CP-44 & Completa\\ \hline
50 & Administrador & \id{RF-37} & CU-07 & M-02 & P-02 & No aplica & CA-21 & CP-45 & Completa\\ \hline
51 & Administrador & \id{RF-38} & CU-07 & M-02 & P-02 & No aplica & CB-17 & CP-46 & Completa\\ \hline
52 & Administrador & \id{RF-39} & CU-07 & M-02 & P-02 & No aplica & CB-18 & CP-47 & Completa\\ \hline
53 & Trabajador & \id{RF-40} & CU-13 & M-03 & P-03 & No aplica & CA-22 & CP-53 & Completa\\ \hline
54 & Trabajador & \id{RF-41} & CU-13 & M-03 & P-03 & No aplica & CA-23 & CP-54 & Completa\\ \hline
55 & Administrador & \id{RF-42} & CU-13 & M-03 & P-03 & No aplica & CA-24 & CP-55 & Completa\\ \hline
56 & Asist. admin. & \id{RF-IA-01} & CU-01 & M-05 & P-05 & A: Abierta $\to$ Notificada & en ficha & CP-56 & Completa\\ \hline
57 & Asist. admin. & \id{RF-IA-02} & CU-01 & M-05 & P-05 & No aplica & en ficha & CP-57 & Completa\\ \hline
58 & Admin. & \id{RF-IA-03} & CU-18 & M-05 & P-05 & A: $\to$ Descartada & en ficha & CP-58 & Completa\\ \hline
59 & Téc. extr. & \id{RF-IA-04} & CU-06 & M-06 & P-06 & R: En desarrollo $\to$ Maduro & en ficha & CP-59 & Completa\\ \hline
60 & Téc. extr. & \id{RF-IA-05} & CU-06 & M-06 & P-06 & R: Maduro $\to$ Cosechado & en ficha & CP-60 & Completa\\ \hline
61 & Admin. & \id{RF-IA-06} & CU-10 & M-06 & P-06 & R: Cosechado $\to$ Despachado & en ficha & CP-61 & Completa\\ \hline
62 & Asist. admin. & \id{RNF-IA-01} & CU-01 & M-05 & P-05$^{*}$ & No aplica & en ficha & CP-62 & Completa\\ \hline
63 & Asist. admin. & \id{RNF-IA-02} & CU-01 & M-05 & P-05$^{*}$ & No aplica & en ficha & CP-63 & Completa\\ \hline
64 & Trabajador & \id{RNF-IA-03} & CU-01 & M-05 & P-05$^{*}$ & No aplica & en ficha & CP-64 & Completa\\ \hline
65 & Admin. & \id{RNF-IA-04} & CU-01 & M-05 & P-05$^{*}$ & No aplica & en ficha & CP-65 & Completa\\ \hline
66 & Trabajador & \id{RNF-IA-05} & CU-01 & M-05 & P-05$^{*}$ & No aplica & en ficha & CP-66 & Completa\\ \hline
67 & Asist. admin. & \id{RNF-IA-06} & CU-18 & M-05 & P-05$^{*}$ & No aplica & en ficha & CP-67 & Completa\\ \hline
68 & Téc. extr. & \id{RNF-IA-07} & CU-06 & M-06 & P-06$^{*}$ & No aplica & en ficha & CP-68 & Completa\\ \hline
69 & Téc. extr. & \id{RNF-IA-08} & CU-06 & M-06 & P-06$^{*}$ & No aplica & en ficha & CP-69 & Completa\\ \hline
70 & Admin. & \id{RNF-IA-09} & CU-06 & M-06 & P-06$^{*}$ & No aplica & en ficha & CP-70 & Completa\\ \hline
71 & Trabajador & \id{RNF-IA-10} & CU-10 & M-06 & P-06$^{*}$ & No aplica & en ficha & CP-71 & Completa\\ \hline
72 & Admin. & \id{RNF-IA-11} & CU-10 & M-06 & P-06$^{*}$ & No aplica & en ficha & CP-72 & Completa\\ \hline
73 & Admin. & \id{RNF-IA-12} & CU-18 & M-06 & P-06$^{*}$ & No aplica & en ficha & CP-73 & Completa\\ \hline
\end{longtable}}

\subsubsection*{Diagnóstico y elementos huérfanos}

\begin{table}[H]
\centering\footnotesize
\caption{Diagnóstico de la trazabilidad}
\label{tab:diag-traza}
\begin{tabular}{|L{7.0cm}|C{2.4cm}|L{6.2cm}|}
\hline
\rowcolor{grisclaro}\textbf{Indicador} & \textbf{Valor} & \textbf{Lectura}\\ \hline
Filas de la matriz & 73 & Cubre 42 RF, 10 RNF, 3 RD y los 18 requisitos de IA\\ \hline
Trazas completas & 73 (100\,\%) & Ninguna fila carece de un eslabón exigible\\ \hline
Trazas parciales & 0 & ---\\ \hline
Trazas huérfanas & 0 & ---\\ \hline
Filas con estado modelado & 26 & Las 47 restantes declaran \emph{No aplica}\\ \hline
Filas con criterio BDD & 71 & Las 2 restantes son restricciones de diseño\\ \hline
Filas sin prototipo & 1 & Requisito de prioridad distinta de \emph{Must}\\ \hline
\end{tabular}
\end{table}

\noindent
La matriz no arroja elementos huérfanos, pero sí tres vacíos que se declaran de forma expresa en
lugar de rellenarlos:

\begin{enumerate}[leftmargin=*, nosep]
  \item \textbf{Requisitos sin prototipo de interfaz.} \id{RF-03}, \id{RF-27} y \id{RF-34}. Los tres
        son de prioridad distinta de \emph{Must}, de modo que el vacío no incumple la exigencia de
        prototipar las pantallas críticas. Se difiere a la Entrega~4.
  \item \textbf{Restricciones sin criterio de aceptación.} \id{RD-07} y \id{RD-10}. Una restricción
        de diseño no describe un comportamiento del sistema sino una condición impuesta sobre su
        construcción, y por tanto no admite criterio de aceptación en el mismo sentido que un
        requisito. Quedan fuera del denominador de la métrica de verificabilidad.
  \item \textbf{Requisitos de IA sin criterio BDD en columna propia.} Los dieciocho requisitos de
        inteligencia artificial llevan su criterio de aceptación dentro de la propia ficha de la \S9
        del ERS, en formato Dado--Cuando--Entonces, y por eso la columna remite a ella en lugar de
        citar un identificador \id{CA-XX} o \id{CB-XX}. La decisión evita duplicar el criterio en dos
        lugares del documento, que es precisamente el mecanismo que produjo los defectos de
        consistencia detectados en la Unidad~IV.
  \item \textbf{Requisitos no funcionales sin proceso propio.} Un requisito no funcional no realiza
        un proceso: lo condiciona. Se les asigna el proceso mayoritario de los requisitos funcionales
        que comparten su caso de uso, y la matriz lo marca con asterisco para que la derivación sea
        visible y no se confunda con una asignación directa.
\end{enumerate}

\subsubsection*{Sincronización entre el ERS y el tablero CASE}

La trazabilidad solo es útil si el tablero refleja el documento. El contraste cuantitativo entre
ambos artefactos, antes y después de la actualización del backlog, es el siguiente:

\begin{table}[H]
\centering\footnotesize
\caption{Sincronización entre el ERS y el tablero de Jira}
\label{tab:sincro}
\begin{tabular}{|L{4.6cm}|C{1.5cm}|C{1.6cm}|C{1.6cm}|C{1.6cm}|C{1.6cm}|}
\hline
\rowcolor{grisclaro}\textbf{Elemento} & \textbf{En el ERS} & \textbf{Tablero inicial} & \textbf{Diverg. inicial} & \textbf{Tablero final} & \textbf{Diverg. final}\\ \hline
Epics / bloques funcionales & 10 & 7 & 3 & 10 & 0\\ \hline
Stories / requisitos & 79 & 39 & 40 & 79 & 0\\ \hline
Sub-tasks / criterios de aceptación & 79 & 38 & 41 & 86 & 0\\ \hline
Historias de usuario & 24 & 19 & 5 & 24 & 0\\ \hline
Filas de trazabilidad & 73 & 52 & 21 & 73 & 0\\ \hline
\multicolumn{3}{|r|}{\textbf{Índice de sincronización}} & \textbf{58,5\,\%} & & \textbf{100\,\%}\\ \hline
\end{tabular}
\end{table}

\noindent
La cifra publicada al cierre de la Unidad~IV fue 67,1\,\%. Se conserva aquí porque es la que consta
en aquel informe, pero el denominador con que se calculó era incompleto: contabilizaba 48 elementos
especificados donde el ERS especifica 79, porque no computaba los diecinueve requisitos no
funcionales ni los doce \id{RNF-IA} como elementos que el tablero debía reflejar. Recalculada sobre
el denominador completo, la sincronización de partida era del 58,5\,\%.

Que el punto de partida fuera peor que el declarado no es un detalle aritmético. Significa que el
indicador que debía medir la desincronización estaba él mismo desincronizado: se ignoraba no solo
cuántos elementos faltaban, sino cuántos debían estar. Un instrumento de medición cuyo denominador
depende de lo que se recuerde incluir mide la memoria del equipo, no el estado del tablero.

La divergencia inicial tiene dos causas verificables. La primera, que la carga del backlog se
ejecutó en la Unidad~IV, antes de que el Change Control Board aprobara \id{RFC-01} a \id{RFC-03}:
los cambios se implementaron en el documento y no se propagaron al tablero. La segunda, que los
requisitos no funcionales nunca llegaron a cargarse, hecho que no se detectó hasta auditar el
tablero elemento por elemento en esta unidad.

El hallazgo es más interesante que el número. Un backlog desincronizado no es un descuido puntual:
es la consecuencia previsible de mantener dos representaciones del mismo conjunto de requisitos sin
un mecanismo que las reconcilie. La misma causa produjo, en la Unidad~IV, once de los treinta y tres
defectos de consistencia detectados en la inspección. La acción correctiva registrada no es solo
cargar los elementos que faltan, sino incorporar al procedimiento de cierre de cada RFC la
actualización del tablero como paso obligatorio, del mismo modo que ya lo es la actualización del
\id{CHANGELOG.md}.

Al cierre de esta unidad la sincronización es del 100\,\%: los 175 elementos del tablero ---diez
epics, setenta y nueve \emph{stories} y ochenta y seis \emph{sub-tasks}--- cubren la totalidad de lo
especificado. Las siete \emph{sub-tasks} que exceden a los criterios del documento son tareas de
verificación creadas por el equipo para su propio seguimiento y no derivan de un criterio de
aceptación, por lo que no se computan como divergencia: el índice mide si algo especificado quedó
sin cargar, no si el tablero contiene elementos adicionales. El respaldo es la exportación
\id{backlog\_export.csv} del Apéndice de trazabilidad.

\input{sec7_ia}
\section{Métricas de calidad del ERS}
\label{sec:metricas}

\subsection{Diseño del instrumento de auditoría}

La auditoría emplea seis métricas derivadas del modelo de calidad de producto de
ISO/IEC 25010:2023 \cite{iso25010} y de las características de un buen conjunto de requisitos de
ISO/IEC/IEEE 29148:2018 \cite{iso29148}. Todas se expresan como una fracción sobre un total contable,
y todos los conteos se obtuvieron por análisis programático de las fuentes \LaTeX{} del ERS.

Esa última decisión merece justificarse. La alternativa ---contar a mano--- es más rápida y produce
cifras que nadie puede reproducir. Contar sobre las fuentes significa que el denominador de cada
métrica es el resultado de un patrón declarado: el número de requisitos funcionales es el número de
macros \verb|\RF{}| al inicio de línea, y el número de casos de uso especificados es el número de
tablas con \verb|\caption{CU-XX}|. Cualquier evaluador puede repetir el conteo y obtener el mismo
número, que es la condición mínima para que una métrica sea una medición y no una opinión con
decimales.

\begin{table}[H]
\centering\footnotesize
\caption{Conteos base de la auditoría}
\label{tab:conteos}
\begin{tabular}{|L{7.2cm}|C{1.8cm}|L{6.0cm}|}
\hline
\rowcolor{grisclaro}\textbf{Concepto} & \textbf{Valor} & \textbf{Procedencia}\\ \hline
Requisitos funcionales & 42 & Macros \verb|\RF{}| al inicio de línea\\ \hline
Requisitos no funcionales & 19 & Filas \id{RNF-XX} de la Tabla 49\\ \hline
Total auditado (RF + RNF) & 61 & Suma de los anteriores\\ \hline
Casos de uso identificados & 18 & Referencias \id{CU-XX} en diagrama y matriz\\ \hline
Casos de uso especificados & 18 & Tablas con \verb|\caption{CU-XX}|\\ \hline
Clases de usuario & 6 & Tabla de clases de usuario de la \S2.3\\ \hline
Requisitos con criterio BDD & 61 & 24 criterios \id{CA-XX} y 37 criterios \id{CB-XX}\\ \hline
Filas de la matriz & 73 & \id{TR-01} a \id{TR-73}\\ \hline
Defectos residuales & 7 & Re-inspección de la \S\ref{sec:validacion}\\ \hline
\end{tabular}
\end{table}

\subsection{Tabla de auditoría con la aritmética de cada métrica}

\begin{longtable}{|L{3.4cm}|L{3.2cm}|C{1.8cm}|C{1.5cm}|C{1.4cm}|C{1.3cm}|}
\hline
\rowcolor{grisclaro}\textbf{Métrica} & \textbf{Fórmula} & \textbf{Aritmética} & \textbf{Valor} & \textbf{Ref.} & \textbf{Cumple}\\ \hline
\endfirsthead
\hline
\rowcolor{grisclaro}\textbf{Métrica} & \textbf{Fórmula} & \textbf{Aritmética} & \textbf{Valor} & \textbf{Ref.} & \textbf{Cumple}\\ \hline
\endhead
M1a Completitud de atributos & requisitos con los 4 atributos / requisitos & $(42{+}19)/61$ & 100\,\% & $\geq 95$\,\% & Sí\\ \hline
M1b Especificación de casos de uso & CU especificados / CU identificados & $18/18$ & 100\,\% & 100\,\% & Sí\\ \hline
M1c Cobertura de actores & actores con $\geq 1$ RF / actores & $6/6$ & 100\,\% & 100\,\% & Sí\\ \hline
M2 Consistencia & $1 -$ conflictos / pares analizados & $1-2/210$ & 0,990 & $\geq 0{,}98$ & Sí\\ \hline
M3 Verificabilidad & requisitos con criterio BDD / requisitos & $61/61$ & 100\,\% & $\geq 90$\,\% & Sí\\ \hline
M3b Verificabilidad en \emph{Must} & \emph{Must} con criterio BDD / \emph{Must} & $24/24$ & 100\,\% & 100\,\% & Sí\\ \hline
M4ade Trazabilidad adelante & RF con cadena completa / RF & $42/42$ & 100\,\% & $\geq 90$\,\% & Sí\\ \hline
M4atr Trazabilidad atrás & RF con fuente / RF & $42/42$ & 100\,\% & 100\,\% & Sí\\ \hline
M5 Modificabilidad & $\sum$ requisitos afectados / muestra & $11/5$ & 2,20 & $\leq 3{,}0$ & Sí\\ \hline
M6 Corrección & defectos residuales / requisitos & $7/61$ & 0,11 & $\leq 0{,}05$ & \textbf{No}\\ \hline
\caption{Auditoría de calidad del ERS final}
\label{tab:auditoria}
\end{longtable}

\subsection{Comparación antes y después de las acciones de mejora}

\begin{table}[H]
\centering\footnotesize
\caption{Métricas antes y después de las acciones de mejora}
\label{tab:antes-despues}
\begin{tabular}{|L{3.6cm}|C{1.6cm}|C{1.6cm}|C{1.5cm}|L{6.0cm}|}
\hline
\rowcolor{grisclaro}\textbf{Métrica} & \textbf{Antes} & \textbf{Después} & \textbf{Ref.} & \textbf{Acción aplicada}\\ \hline
M1a Completitud & 98,4\,\% & 100\,\% & $\geq 95$\,\% & \id{RES-04}: se clasificó \id{RNF-19} en el Apéndice C\\ \hline
\rowcolor{grissuave}M1b Casos de uso & \textbf{55,6\,\%} & 100\,\% & 100\,\% & \id{RES-05}: se especificaron \id{CU-11} a \id{CU-18}\\ \hline
\rowcolor{grissuave}M1c Actores & \textbf{75,0\,\%} & 100\,\% & 100\,\% & \id{RES-07}: reconteo tras distinguir clase de usuario de parte interesada\\ \hline
M2 Consistencia & 0,990 & 0,990 & $\geq 0{,}98$ & Sin acción; los dos conflictos quedaron cerrados\\ \hline
\rowcolor{grissuave}M3 Verificabilidad & \textbf{39,3\,\%} & 100\,\% & $\geq 90$\,\% & \id{RES-06}: 37 criterios reescritos en formato BDD\\ \hline
M4 Trazabilidad & 100\,\% & 100\,\% & $\geq 90$\,\% & Los DFD aportan el eslabón de proceso\\ \hline
M5 Modificabilidad & 2,20 & 2,20 & $\leq 3{,}0$ & Sin acción\\ \hline
\rowcolor{rojosuave}M6 Corrección & \textbf{0,11} & \textbf{0,11} & $\leq 0{,}05$ & Sin acción posible; véase la lectura siguiente\\ \hline
\end{tabular}
\end{table}

\subsection{Lectura interpretativa de cada métrica}

\subsubsection*{M1 Completitud: lo que una métrica ve y una lectura no}

Los tres indicadores de completitud arrojaron resultados muy distintos: 98,4\,\% en atributos,
55,6\,\% en casos de uso y 75,0\,\% en cobertura de actores. Que el primero fuera casi perfecto y el
segundo estuviera a la mitad dice algo sobre cómo se degrada un documento: los elementos que existen
están bien especificados, y lo que falta son elementos enteros que nadie echa en falta porque no
aparecen en ninguna parte salvo en un diagrama.

Ninguna lectura del documento habría detectado que faltaban ocho casos de uso. La métrica lo detectó
en una división.

\subsubsection*{M2 Consistencia: el denominador es la decisión}

M2 se calcula como uno menos la razón entre conflictos detectados y pares analizados. La cifra
depende enteramente del denominador, y ahí está la decisión de método: los sesenta y un requisitos
producen 1.830 pares posibles, y compararlos todos habría dado un denominador que hace tender la
métrica a uno por construcción.

Se analizaron 210 pares: los que pertenecen al mismo bloque funcional o comparten almacén de datos,
que son los que pueden contradecirse. Declarar el denominador y su criterio es lo que hace la métrica
interpretable; un 0,990 sobre 1.830 pares y un 0,990 sobre 210 no son la misma afirmación.

\subsubsection*{M3 Verificabilidad: qué cambió realmente}

La métrica pasó de 39,3\,\% a 100\,\%, y el salto es el mayor de la auditoría. Conviene precisar qué
lo produjo para no sobreinterpretarlo: no se especificó ningún requisito nuevo ni se modificó ningún
umbral. Los treinta y siete criterios reescritos conservan el valor, la unidad y el método que su
ficha ya declaraba.

Lo que cambió es la forma de enunciar la condición, y con ella la posibilidad de que dos evaluadores
independientes lleguen al mismo veredicto. Un ERS con el cien por cien de sus requisitos en formato
BDD no es un ERS correcto: es un ERS cuyos requisitos son comprobables. La corrección se mide aparte,
con M6, y su valor es peor.

\subsubsection*{M4 Trazabilidad: la métrica que obligó a producir un modelo}

M4 hacia adelante exige que cada requisito enlace con caso de uso, módulo, proceso y prueba. El
eslabón de proceso no existía porque el ERS no tenía diagramas de flujo de datos, de modo que la
métrica no podía calcularse sin producir antes el modelo que faltaba.

Es el único caso de la auditoría en que medir obligó a construir. Y el resultado indirecto fue más
valioso que el directo: los DFD no solo cerraron la columna, sino que hicieron explícita la
correspondencia entre bloque funcional, módulo y proceso, que hasta entonces era una convención
tácita del equipo.

\subsubsection*{M5 Modificabilidad: por qué 2,20 es un buen número}

El acoplamiento medio de 2,20 significa que modificar un requisito obliga a revisar algo más de dos
requisitos adicionales. El valor se mantiene bajo por una razón concreta y no por suerte: la
correspondencia uno a uno entre bloque funcional, módulo y proceso del DFD hace que un cambio se
propague por una sola ruta.

La muestra de cinco requisitos combina bloques distintos y las tres prioridades MoSCoW, para no medir
solo el núcleo del sistema. El requisito de mayor acoplamiento resultó ser \id{RF-28}, el registro de
avance, del que dependen la liquidación, el indicador de desempeño y la rectificación: es
consistente con que sea también el requisito con más casos de uso asociados.

\subsubsection*{M6 Corrección: la única métrica que no puede corregirse}

M6 es la única que no cumple, y es la única cuya corrección no depende de actuar sobre el documento.
Mide qué proporción de defectos escapó a la inspección de la Unidad IV: siete defectos residuales
sobre sesenta y un requisitos, es decir 0,11 frente al 0,05 de referencia.

Los siete están corregidos. Corregirlos, sin embargo, no altera el numerador, porque el numerador es
un hecho sobre la eficacia de aquella inspección y no sobre el estado actual del documento. Bajar la
cifra exigiría declarar una re-inspección que no se hizo o descontar los defectos ya reparados, y
cualquiera de las dos cosas convertiría la métrica en un adorno.

Se reporta tal cual. La acción correctiva registrada es de proceso y no de producto: la lista de
verificación de la inspección debe incluir los apéndices y los conteos globales, y debe incorporar
preguntas sobre la completitud del conjunto y no solo sobre la corrección de sus elementos. Es donde
se concentraron seis de los siete.

\subsection{Valoración global}

\begin{table}[H]
\centering\footnotesize
\caption{Valoración ponderada de la auditoría}
\label{tab:valoracion}
\begin{tabular}{|L{4.6cm}|C{1.6cm}|C{1.8cm}|C{1.6cm}|L{4.6cm}|}
\hline
\rowcolor{grisclaro}\textbf{Métrica} & \textbf{Peso} & \textbf{Nivel (0--4)} & \textbf{Aporte} & \textbf{Criterio del nivel}\\ \hline
M1 Completitud & 0,20 & 4 & 0,80 & Cumple con evidencia completa tras la acción correctiva\\ \hline
M2 Consistencia & 0,15 & 4 & 0,60 & Cumple; conflictos detectados, documentados y cerrados\\ \hline
M3 Verificabilidad & 0,20 & 4 & 0,80 & Cumple; 61 de 61 requisitos con criterio comprobable\\ \hline
M4 Trazabilidad & 0,20 & 4 & 0,80 & Cumple en ambos sentidos, sin huérfanos\\ \hline
M5 Modificabilidad & 0,10 & 4 & 0,40 & Cumple sobre muestra representativa\\ \hline
M6 Corrección & 0,15 & 2 & 0,30 & Por debajo de la referencia, con acción correctiva de proceso aplicada\\ \hline
\textbf{Total} & \textbf{1,00} & --- & \textbf{3,70} & Sobre 4,00\\ \hline
\end{tabular}
\end{table}

\noindent
El 3,70 sobre 4,00 se sostiene en cinco métricas cumplidas y una incumplida con acción documentada.
La lectura honesta del conjunto no es que el ERS sea bueno, sino que el ERS \emph{es ahora} medible:
antes de esta unidad no existía un instrumento capaz de decir si lo era.

\section{Retrospectiva del equipo}
\label{sec:retro}

\subsection{Método}

La retrospectiva se ejecuta con el formato Start--Stop--Continue: qué empezar a hacer, qué dejar de
hacer y qué mantener. El formato obliga a que cada elemento sea accionable y tenga responsable, lo
que lo distingue de una valoración general del trabajo.

Los elementos que siguen son los que el equipo identificó en la sesión de cierre, y cada uno se apoya
en un hecho verificable del proyecto y no en una impresión.

Son doce: cuatro de \emph{Start}, cuatro de \emph{Stop} y cuatro de \emph{Continue}.

\subsection{Start --- qué empezar a hacer}

\begin{table}[H]
\centering\footnotesize
\caption{Retrospectiva: elementos Start}
\label{tab:start}
\begin{tabular}{|C{1.2cm}|L{6.4cm}|L{4.4cm}|C{2.4cm}|}
\hline
\rowcolor{grisclaro}\textbf{N.º} & \textbf{Acción} & \textbf{Hecho que la motiva} & \textbf{Responsable}\\ \hline
S1 & Incluir la actualización del tablero CASE en el procedimiento de cierre de toda RFC, del mismo modo que ya lo está la actualización del \id{CHANGELOG.md} & El índice de sincronización entre el ERS y Jira era del 58,5\,\% ---declarado como 67,1\,\% sobre un denominador incompleto---: 107 elementos especificados no estaban en el tablero, incluidos los diecinueve RNF, que nunca llegaron a cargarse. Cerrado al 100\,\% en esta unidad & Alcívar A.\\ \hline
S2 & Ampliar el alcance de la inspección a los apéndices y a los conteos globales cuando el documento haya sufrido cambios aprobados & Cuatro de los siete defectos residuales se concentran ahí, fuera del alcance declarado en la Unidad IV & Rizzo E.\\ \hline
S3 & Incorporar a la lista de verificación preguntas sobre la completitud del conjunto y no solo sobre la corrección de sus elementos & Ningún inspector detectó que faltaban ocho casos de uso, porque ninguna pregunta obligaba a contar cuántos debían existir & Rizzo E.\\ \hline
S4 & Medir antes de corregir, y no al revés & La auditoría detectó en una división lo que tres lecturas del documento no habían detectado & Macías J.\\ \hline
\end{tabular}
\end{table}

\subsection{Stop --- qué dejar de hacer}

\begin{table}[H]
\centering\footnotesize
\caption{Retrospectiva: elementos Stop}
\label{tab:stop}
\begin{tabular}{|C{1.2cm}|L{6.4cm}|L{4.4cm}|C{2.4cm}|}
\hline
\rowcolor{grisclaro}\textbf{N.º} & \textbf{Acción} & \textbf{Hecho que la motiva} & \textbf{Responsable}\\ \hline
T1 & Dejar de declarar cifras agregadas en varios lugares del documento sin una fuente única & Once de los treinta y tres defectos de la inspección son de consistencia entre secciones que declaran la misma cifra & Huilcapi D.\\ \hline
T2 & Dejar de dar por válida una declaración de conformidad normativa sin verificar la estructura de la norma citada & La cobertura declarada de las nueve características de ISO/IEC 25010:2023 no se sostenía: se empleaba una subcaracterística como característica de primer nivel & Rizzo E.\\ \hline
T3 & Dejar de asumir supuestos sobre el contexto sin dato que los respalde & Dos restricciones de diseño se justificaban con suposiciones que la evidencia \id{EV-12} había invalidado & Villafuerte A.\\ \hline
T4 & Dejar de aceptar el mensaje que la interfaz web de GitHub propone por defecto al subir archivos & 45 de los 84 \emph{commits} del repositorio de la Unidad IV (\id{ISR401-PFC-ERS-EQUIPO\_B}) usan ``Add files via upload'' o ``Rename\ldots'': el historial documenta qué archivo cambió, pero no por qué & Alcívar A.\\ \hline
\end{tabular}
\end{table}

\subsection{Continue --- qué mantener}

\begin{table}[H]
\centering\footnotesize
\caption{Retrospectiva: elementos Continue}
\label{tab:continue}
\begin{tabular}{|C{1.2cm}|L{6.4cm}|L{4.4cm}|C{2.4cm}|}
\hline
\rowcolor{grisclaro}\textbf{N.º} & \textbf{Acción} & \textbf{Hecho que la motiva} & \textbf{Responsable}\\ \hline
C1 & Mantener la preparación individual previa a la inspección, con marca temporal verificable & Los cuatro roles superaron el mínimo de seis defectos y ninguno concentró más del 33\,\% del total: el reparto por secciones funcionó & Macías J.\\ \hline
C2 & Mantener las fuentes del ERS versionadas junto al PDF & Permitió reconstruir, corregir y recompilar el documento en cada unidad sin rehacerlo & Arboleda F.\\ \hline
C3 & Mantener la derivación del análisis de impacto desde la matriz y no por estimación & Las tres RFC declararon identificadores de fila verificables en el artefacto & Arboleda F.\\ \hline
C4 & Mantener la declaración expresa de los vacíos en lugar de rellenarlos & Los requisitos sin prototipo, las restricciones sin criterio BDD y las filas sin estado se declaran; un vacío declarado es un dato, uno oculto es un defecto & Villafuerte A.\\ \hline
\end{tabular}
\end{table}


% =====================================================================
\section{Conclusiones}
\label{sec:conclusiones}

\subsection*{Conclusión 1 --- Unidad I: el contexto determinó el alcance más que la técnica}

La planificación de la elicitación fijó como objetivo cubrir a los seis perfiles de usuario, y esa
decisión ---tomada antes de conocer el dominio--- resultó ser la que más restricciones generó. El
cuestionario ampliado \id{EV-12}, con $n=62$, reveló que el 11,3\,\% del personal no dispone de
teléfono inteligente y que el 74,2\,\% no tiene señal permanente en el lote. Ambos datos
contradijeron suposiciones que el equipo había dado por válidas y forzaron dos requisitos que no
estaban previstos: el registro delegado \id{RF-35} y la operación sin conexión \id{RD-02}.

La conclusión operativa es que la cobertura de la elicitación no es una cuestión de exhaustividad
metodológica sino de alcance del sistema: si el instrumento no hubiera llegado al personal de campo,
el ERS habría especificado un sistema que el 11,3\,\% de sus usuarios no podría utilizar, y el
defecto no habría aparecido hasta el despliegue.

\subsection*{Conclusión 2 --- Unidad II: la evidencia tardía es la más cara}

La evidencia \id{EV-13} ---la hoja de liquidación semanal--- se obtuvo después de haber cerrado la
especificación, y obligó a incorporar cuatro requisitos funcionales nuevos sobre un proceso que la
organización ejecuta todas las semanas y que consume una jornada administrativa completa. Ninguna de
las doce evidencias anteriores lo había identificado.

El costo de esa incorporación tardía es medible y se distribuyó por todo el documento: obligó a
modificar recuentos en cuatro archivos, produjo la contradicción entre 21 y 19 requisitos
\emph{Must} que la inspección detectó como defecto crítico, y dejó dos requisitos \emph{Must} sin
historia de usuario que hubo que tramitar ante el Change Control Board. Un requisito que llega tarde
no cuesta lo que cuesta especificarlo: cuesta lo que cuesta reconciliar el documento entero con él.

\subsection*{Conclusión 3 --- Unidad III: dos taxonomías del mismo sistema producen defectos}

Once de los treinta y tres defectos detectados en la inspección son de consistencia, y ninguno es de
redacción. El patrón tiene una causa identificable: el documento mantenía la parte documental
---bloques funcionales, fichas de requisito--- y la parte de priorización ---MoSCoW, WSJF, matriz---
como dos representaciones del mismo conjunto de requisitos, sin mecanismo que las reconciliara.

La Unidad V confirmó el diagnóstico desde otro ángulo. Al incorporar los diagramas de flujo de datos,
la decisión deliberada fue hacer corresponder uno a uno los siete procesos con los siete bloques
funcionales existentes, en lugar de descomponer el sistema por criterio propio del análisis de flujo.
El modelo resultante es posiblemente menos fiel que el que habría producido una descomposición
independiente, y se aceptó esa pérdida precisamente para no crear una tercera taxonomía paralela. El
acoplamiento medio de 2,20 de la métrica M5 es la consecuencia medible de esa decisión.

\subsection*{Conclusión 4 --- Unidad IV: la inspección es barata y el alcance que se le fija es lo caro}

La inspección formal detectó treinta y tres defectos a un costo de 0,44 horas-persona por defecto, lo
que sostiene empíricamente el argumento de Fagan sobre el costo relativo de la detección temprana
\cite{fagan1976}. Pero la re-inspección de la Unidad V encontró siete defectos residuales, y su
distribución es más informativa que su número: cuatro se concentran en los apéndices y en los
conteos globales, que el alcance de la inspección había excluido de forma expresa.

La exclusión era razonable cuando se decidió ---los apéndices son material de referencia--- y dejó de
serlo en cuanto tres solicitudes de cambio modificaron los recuentos que allí residen. Nadie revisó
la decisión de alcance después de aprobar las RFC. La conclusión no es que la inspección fallara,
sino que un alcance de inspección es una decisión con fecha de caducidad: caduca cuando cambia el
documento sobre el que se fijó.

\subsection*{Conclusión 5 --- Unidad V: especificar equidad exige preguntar entre quiénes}

El ERS de la Entrega 2A especificaba exactitud, latencia y explicabilidad para sus componentes
automáticos, y no especificaba equidad. La ausencia no fue un descuido de redacción: fue un punto
ciego del modelo de calidad empleado, que no contemplaba la propiedad, y solo se hizo visible al
formular explícitamente la pregunta de entre qué grupos podía diferir el desempeño del sistema.

La respuesta, en un cultivo de palma, no es demográfica. Los cuatro requisitos de equidad
especificados comparan variedades sembradas, gamas de dispositivo y cuadrillas de trabajo. El de
mayor consecuencia es \id{RNF-IA-10}: como \id{RF-24} traza la penalización de la extractora hasta la
cuadrilla responsable, un modelo que estimara peor las cargas de una cuadrilla concreta produciría
recortes de remuneración sistemáticamente injustos sobre personas identificables. Es el único
indicador del plan de monitoreo cuya acción al superarse es suspender el uso del modelo y no
reentrenarlo.

La conclusión que integra las cinco unidades es esta: el proyecto midió lo que sabía nombrar. Los
defectos de consistencia aparecieron cuando se contaron las cifras; los casos de uso faltantes,
cuando se dividió especificados entre identificados; y la ausencia de requisitos de equidad, cuando
se preguntó entre qué grupos. En los tres casos el problema existía antes de la medición y llevaba
meses en el documento sin que ninguna lectura lo revelara. La ingeniería de requisitos no consiste
solo en escribir bien lo que se sabe, sino en construir los instrumentos que hacen visible lo que
todavía no se ha nombrado.


%=======================================================================
% REFERENCIAS
%=======================================================================
\newpage
\IfFileExists{IEEEtran.bst}{\bibliographystyle{IEEEtran}}{\bibliographystyle{unsrt}}
\bibliography{referencias}
\addcontentsline{toc}{section}{Referencias}

%=======================================================================
% ANEXOS
%=======================================================================
\newpage
\appendix
\section*{Anexos}
\addcontentsline{toc}{section}{Anexos}
\setcounter{section}{0}
\renewcommand{\thesection}{\Alph{section}}

\section{Instrumento de auditoría de calidad}

El instrumento completo se versiona en \id{02\_Auditoria/AnexoA\_auditoria\_calidad.xlsx} y contiene
seis hojas: \emph{Auditoria} con la tabla de las seis métricas y sus fórmulas vivas;
\emph{Antes\_Despues} con la comparación de las dos mediciones; \emph{Conteos\_base} con el valor y la
procedencia de cada conteo; \emph{Conflictos\_M2} con los pares analizados y los conflictos hallados;
\emph{Trazabilidad\_M4} con el recuento por eslabón; \emph{Modificabilidad\_M5} con la muestra de
acoplamiento; y \emph{Reinspeccion\_M6} con los defectos residuales.

Ninguna celda de valor está escrita a mano: todas son fórmulas sobre la hoja de conteos, de modo que
modificar un conteo recalcula la métrica y su veredicto. La tabla resumen se reproduce en la
\S\ref{sec:metricas}.

\section{Banco de preguntas del tribunal y respuestas anticipadas}

Las veintidós preguntas siguientes son las del Anexo B de la guía. Para cada una se registra la
respuesta preparada por el equipo y el artefacto donde se verifica, conforme al Paso 4.b.

\subsection*{Fundamentos}

\begin{longtable}{|C{0.7cm}|L{14.6cm}|}
\hline
\rowcolor{grisclaro}\textbf{1} & \textbf{¿Cuál es la frontera del sistema y qué queda deliberadamente fuera?}\\ \hline
& La frontera se declara en el DFD de nivel 0 (\S8.2 del ERS). Dentro quedan la captura de datos en campo, el diagnóstico, la planificación, la liquidación y los reportes. Fuera quedan tres cosas por decisión y no por omisión: la contabilidad y la nómina formal de la organización, que se llevan en otro sistema; la operación de la planta extractora, de la que solo se recibe la calificación de calidad; y la ejecución de la aplicación fitosanitaria, sobre la que el sistema recomienda pero nunca decide. \newline \emph{Artefacto:} \S8.2 del ERS, Figura del DFD de nivel 0.\\ \hline
\rowcolor{grisclaro}\textbf{2} & \textbf{¿Qué diferencia hay entre el contexto y el límite del sistema en su proyecto?}\\ \hline
& El contexto es todo lo que influye en el sistema aunque no interactúe con él: la estacionalidad del cultivo, la política de penalización de la extractora, la LOPDP. El límite es la línea que separa lo que el sistema hace de lo que hacen otros. La distinción tuvo consecuencia práctica: la asesoría técnica es contexto ---aporta fichas y recibe diagnósticos--- pero no es clase de usuario, porque no opera la aplicación. Confundir ambas cosas fue lo que produjo el defecto \id{RES-07} en la primera medición de la cobertura de actores. \newline \emph{Artefacto:} \S2.3 y \S2.4 del ERS; recuadro sobre clase de usuario y parte interesada.\\ \hline
\rowcolor{grisclaro}\textbf{3} & \textbf{¿Qué modelo de proceso siguieron y por qué era el adecuado para su dominio?}\\ \hline
& Un proceso documental incremental, no ágil puro. La razón es que ocho requisitos legales de la LOPDP deben poder demostrarse ante un tercero, y una obligación legal no se satisface con una historia priorizada. Pero se mantuvieron historias de usuario y criterios en formato BDD porque el dominio es volátil: la evidencia \id{EV-13} apareció después de cerrar la especificación y añadió cuatro requisitos. Convivir con ambas representaciones tuvo un costo medible, y lo declaramos: once de los treinta y tres defectos de la inspección son de consistencia entre ellas. \newline \emph{Artefacto:} \S2.3 del informe; \S\ref{sec:metricas}, lectura de M2.\\ \hline
\end{longtable}

\subsection*{Elicitación}

\begin{longtable}{|C{0.7cm}|L{14.6cm}|}
\hline
\rowcolor{grisclaro}\textbf{4} & \textbf{¿Qué técnica les aportó los requisitos que no habrían descubierto de otro modo?}\\ \hline
& El cuestionario ampliado \id{EV-12}, con $n=62$. Las entrevistas con la administración describieron los procesos formales; el cuestionario al personal de campo reveló que el 11,3\,\% no dispone de teléfono inteligente y que el 74,2\,\% no tiene señal permanente en el lote. Ninguno de los dos datos habría salido de una entrevista con la administración, porque la administración no los conocía. De ahí nacieron \id{RF-35} y \id{RD-02}. \newline \emph{Artefacto:} \id{EV-12}; \S2.6 y \S2.7 del ERS.\\ \hline
\rowcolor{grisclaro}\textbf{5} & \textbf{¿Qué \emph{stakeholder} quedó subrepresentado y cómo lo compensaron?}\\ \hline
& El trabajador agrícola. Es la clase de usuario más numerosa y la de competencia digital declarada como baja, y el formato de entrevista no funcionó con ella. Se compensó con dos instrumentos: el cuestionario \id{EV-12}, que alcanzó a las sesenta y dos personas, y la observación directa en campo. Aun así la compensación es parcial y se declara como tal: el trabajador aportó datos sobre su contexto, no requisitos formulados por él. \newline \emph{Artefacto:} \S2.3 del ERS; \id{EV-05}, \id{EV-06}, \id{EV-12}.\\ \hline
\rowcolor{grisclaro}\textbf{6} & \textbf{Muéstrenme un requisito y díganme de qué evidencia concreta nació.}\\ \hline
& \id{RF-37}, cálculo de la remuneración semanal. Nació de \id{EV-13}, la hoja de liquidación que la organización facilitó y que documenta que el cálculo manual consume una jornada administrativa completa cada semana. La evidencia llegó después de cerrar la especificación, lo que obligó a incorporar \id{RF-36} a \id{RF-39} y produjo en cascada el defecto crítico \id{DEF-02}. \newline \emph{Artefacto:} \id{EV-13}; ficha de \id{RF-37} en la \S3.2 del ERS.\\ \hline
\end{longtable}

\subsection*{Especificación}

\begin{longtable}{|C{0.7cm}|L{14.6cm}|}
\hline
\rowcolor{grisclaro}\textbf{7} & \textbf{Elijan un RNF y expliquen cómo se comprueba objetivamente.}\\ \hline
& \id{RNF-10}, disponibilidad mensual del servicio de respaldo igual o superior al 99\,\%. Se comprueba totalizando el tiempo de indisponibilidad registrado en un mes natural y verificando que no supera 7~h~12~min, que es el 1\,\% de las 720~h del período. El criterio \id{CB-28} lo enuncia en formato Dado--Cuando--Entonces. El valor de 7~h~12~min es, además, una corrección: el ERS declaraba 7~h~18~min y la inspección lo detectó como defecto \id{DEF-20}. \newline \emph{Artefacto:} Tabla 49 del ERS; criterio \id{CB-28} del Apéndice F.\\ \hline
\rowcolor{grisclaro}\textbf{8} & \textbf{¿Qué requisito les costó más redactar de forma verificable y por qué?}\\ \hline
& \id{RF-16}, el indicador de desempeño del personal. La versión original decía que el sistema calcula un indicador, y su criterio de verificación era que el sistema calcula un indicador: cualquier salida lo satisfacía. La inspección lo registró como \id{DEF-18}. Corregirlo exigió fijar fórmula, tres componentes, sus pesos y una escala de 0 a 100, y esos valores los eligió el equipo: están declarados como decisión de análisis pendiente de validar con la organización. \newline \emph{Artefacto:} \S5.5 del informe, valores fijados durante la corrección.\\ \hline
\rowcolor{grisclaro}\textbf{9} & \textbf{¿Por qué este caso de uso incluye a aquel y no lo extiende?}\\ \hline
& \id{CU-01}, \id{CU-02} y \id{CU-06} \textbf{incluyen} a \id{CU-18} porque la explicación de una decisión automática ocurre siempre, no bajo condición: sin ella el caso de uso base no está completo. En cambio \id{CU-15} \textbf{extiende} con \id{CU-16} porque la alerta solo se genera si el porcentaje de afectación supera el umbral. La regla que aplicamos es que la inclusión es obligatoria e incondicional y la extensión es condicional. \newline \emph{Artefacto:} \S4 del ERS, relaciones de inclusión y extensión.\\ \hline
\end{longtable}

\subsection*{Validación}

\begin{longtable}{|C{0.7cm}|L{14.6cm}|}
\hline
\rowcolor{grisclaro}\textbf{10} & \textbf{¿Cuántos defectos halló la inspección y cuáles siguen abiertos?}\\ \hline
& Treinta y tres defectos únicos: 4 críticos, 24 mayores y 5 menores. Se corrigieron los veintiocho críticos y mayores ---veinticinco directamente y tres mediante RFC--- y cuatro de los cinco menores. Sigue abierto uno solo: \id{DEF-25}, sobre \id{RF-09}, que exige recomendar la acción de control según la etapa del insecto sin que las etapas estén enumeradas. Corregirlo requiere el ciclo biológico de cada plaga, información entomológica que ninguna de las trece evidencias aporta. Preferimos dejarlo declaradamente incompleto antes que enumerar etapas sin sustento. \newline \emph{Artefacto:} Tabla \ref{tab:defectos}; \S5 del informe PE4.\\ \hline
\rowcolor{grisclaro}\textbf{11} & \textbf{¿Qué defecto crítico se les escapó y cómo lo detectaron después?}\\ \hline
& Ninguno de severidad crítica, pero sí seis mayores y uno menor, detectados en la re-inspección de esta unidad. Cuatro se concentran en los apéndices y en los conteos globales, que el alcance de la inspección había excluido de forma expresa; dos son de completitud del conjunto ---faltaban ocho casos de uso y treinta y siete criterios BDD--- y no de corrección de sus elementos. Los detectó la auditoría cuantitativa, no una lectura: M1b arrojó 55,6\,\% al dividir casos especificados entre identificados. \newline \emph{Artefacto:} \S\ref{sec:validacion}, análisis de causa; hoja \emph{Reinspeccion\_M6} del Anexo A.\\ \hline
\rowcolor{grisclaro}\textbf{12} & \textbf{¿Cómo saben que las correcciones no introdujeron nuevos defectos?}\\ \hline
& No lo sabemos con certeza, y conviene decirlo así. Lo que hicimos son tres comprobaciones: recompilar el ERS tras cada tanda de correcciones verificando que no aparecen errores ni referencias sin resolver; ejecutar la re-inspección sobre el documento ya corregido, que es precisamente lo que mide M6; y recalcular las seis métricas después de cada acción. Las tres detectan clases distintas de defecto y ninguna las cubre todas. La medición de M6 arrojó 0,11, por encima del 0,05 de referencia, y ese número dice justamente que el proceso de corrección anterior no fue completo. \newline \emph{Artefacto:} \S\ref{sec:metricas}, lectura de M6.\\ \hline
\end{longtable}

\subsection*{Gestión}

\begin{longtable}{|C{0.7cm}|L{14.6cm}|}
\hline
\rowcolor{grisclaro}\textbf{13} & \textbf{Tomen un RF y díganme qué se rompería si mañana cambia.}\\ \hline
& \id{RF-28}, el registro de avance por unidad de labor. Es el requisito de mayor acoplamiento de la muestra de M5: afecta a cuatro requisitos. Se rompen \id{RF-29}, porque el acumulado que consulta el trabajador se calcula sobre él; \id{RF-37}, porque la liquidación semanal multiplica ese avance por la tarifa; \id{RF-16}, porque el indicador de desempeño toma el avance como uno de sus tres componentes; y \id{RF-41}, porque la rectificación opera sobre esos registros. Los cuatro se obtienen filtrando la matriz, no estimando. \newline \emph{Artefacto:} hoja \emph{Modificabilidad\_M5} del Anexo A; matriz de trazabilidad.\\ \hline
\rowcolor{grisclaro}\textbf{14} & \textbf{¿Qué RFC rechazó el CCB y con qué argumento?}\\ \hline
& \textbf{Ninguna fue rechazada}: las tres se aprobaron, una de ellas con condiciones. Conviene explicar por qué, porque un CCB que aprueba todo puede parecer un trámite. Las tres solicitudes no se eligieron de un catálogo de escenarios: son los tres defectos de la inspección cuya reparación consistía en añadir alcance y no en enmendar un texto, de modo que rechazarlas habría significado dejar abierto un defecto crítico o mayor. Lo que sí hubo fue una alternativa rechazada dentro de \id{RFC-03}: se propuso diferir íntegramente la especificación de los requisitos de la LOPDP a la Entrega 4, y se descartó porque habría dejado \id{DEF-12} sin cerrar y habría obligado a rediseñar más tarde tres funciones que tocan el modelo de datos de personal. Se aprobó una fórmula intermedia: especificar ahora, construir después. \newline \emph{Artefacto:} Acta del CCB, deliberación de \id{RFC-03}.\\ \hline
\rowcolor{grisclaro}\textbf{15} & \textbf{¿Cuál es la línea base vigente y dónde está etiquetada?}\\ \hline
& La versión v1.1 del ERS, etiquetada en Git como \id{baseline-v1.1} y, al cierre de esta unidad, como \id{baseline-final}. La versión se declara de forma idéntica en cuatro lugares: portada del ERS, historial de revisiones, \id{CHANGELOG.md} y etiqueta anotada. La comprobación es \verb|git ls-remote --tags origin|, que debe listarla, y la pestaña de \emph{tags} del repositorio. \newline \emph{Artefacto:} \S\ref{sec:gestion}, tabla de versiones; \id{README.md}.\\ \hline
\end{longtable}

\subsection*{Inteligencia artificial}

\begin{longtable}{|C{0.7cm}|L{14.6cm}|}
\hline
\rowcolor{grisclaro}\textbf{16} & \textbf{¿Qué decide el modelo y qué pasa cuando se equivoca?}\\ \hline
& \id{IA-01} clasifica una fotografía en una clase de afectación o en \emph{sin afectación}. Un falso negativo retrasa el control hasta la siguiente visita técnica, que es exactamente la situación actual: el componente no empeora la línea base. Un falso positivo provoca una aplicación innecesaria de producto, con costo y exposición del personal. Por eso el umbral de exhaustividad (0,85) es más exigente que el de precisión. Si la confianza baja del umbral, el componente \emph{no emite clasificación}: declara el resultado no concluyente. \newline \emph{Artefacto:} \S\ref{sec:ia}, bloque 1 de \id{IA-01}; \id{RF-IA-02}.\\ \hline
\rowcolor{grisclaro}\textbf{17} & \textbf{¿Con qué datos lo entrenarían y con qué base legal?}\\ \hline
& Con fotografías capturadas en la propia plantación durante dos ciclos de cosecha, etiquetadas por la asesoría técnica, más conjuntos públicos de fitopatología cuando haya correspondencia de clase. La imagen de una palma no es dato personal; sí lo son la geolocalización y la identidad de quien captura, que se asocian a ella por \id{RF-05}. La base legal es la ejecución de la relación laboral, Art.~7 de la LOPDP, y \textbf{no el consentimiento}: la asimetría de poder propia de la relación laboral haría que un consentimiento otorgado en ese contexto difícilmente fuera libre. Ambos campos se disocian antes de incorporar la imagen al conjunto. \newline \emph{Artefacto:} \S\ref{sec:ia}, bloque 2 y tabla de tratamiento de datos personales.\\ \hline
\rowcolor{grisclaro}\textbf{18} & \textbf{¿Cómo medirían que el modelo no perjudica sistemáticamente a un grupo?}\\ \hline
& Con cuatro requisitos de equidad, y los grupos no son demográficos porque en este dominio el daño no llega por esa vía. \id{RNF-IA-04} y \id{RNF-IA-11} comparan las dos variedades sembradas. \id{RNF-IA-05} compara gamas de dispositivo: como \id{RD-01} establece que el personal usa su propio teléfono, un modelo que rinda peor con cámaras baratas penalizaría a quien tiene el teléfono más económico. El de mayor consecuencia es \id{RNF-IA-10}, que compara cuadrillas con una brecha máxima de 1,5 puntos: \id{RF-24} traza la penalización de la extractora hasta la cuadrilla, de modo que un sesgo produce recortes de remuneración sobre personas identificables. Es el único indicador cuya acción al superarse es \emph{suspender} el uso del modelo. La medición es \emph{a posteriori} sobre el registro de decisiones, porque la cuadrilla se excluye de las variables de entrada. \newline \emph{Artefacto:} \S\ref{sec:ia}, requisitos de equidad y plan de monitoreo.\\ \hline
\end{longtable}

\subsection*{Métricas}

\begin{longtable}{|C{0.7cm}|L{14.6cm}|}
\hline
\rowcolor{grisclaro}\textbf{19} & \textbf{¿Qué métrica salió peor y qué hicieron al respecto?}\\ \hline
& Dos respuestas, porque la pregunta admite dos lecturas. La que salió peor en la primera medición fue M3, verificabilidad, con 39,3\,\% frente al 90\,\% de referencia; se cerró reescribiendo treinta y siete criterios en formato Dado--Cuando--Entonces, y conviene precisar que no se modificó ningún umbral: cambió la forma de enunciar la condición, no su contenido. La que sigue incumpliendo es M6, corrección, con 0,11 frente a 0,05. No la corregimos y no podemos: mide qué proporción de defectos escapó a la inspección anterior, y ese numerador es un hecho histórico. Corregir los siete defectos mejora el documento pero no cambia el número. La acción registrada es de proceso: la lista de verificación debe incluir apéndices y conteos globales. \newline \emph{Artefacto:} \S\ref{sec:metricas}, tabla antes/después y lectura de M6.\\ \hline
\rowcolor{grisclaro}\textbf{20} & \textbf{Enséñenme los conteos con los que calcularon la verificabilidad.}\\ \hline
& M3 $= 61/61 \times 100 =$ 100\,\%. El denominador son 42 requisitos funcionales más 19 no funcionales; las diez restricciones de diseño y los ocho requisitos legales se auditan aparte porque no son requisitos del sistema sino condiciones impuestas. El numerador son 24 criterios \id{CA-01} a \id{CA-24}, procedentes de las historias de usuario, más 37 criterios \id{CB-01} a \id{CB-37} del Apéndice F. Los conteos no se hicieron a mano: el número de requisitos funcionales es el número de macros \verb|\RF{}| al inicio de línea en las fuentes, y cualquiera puede repetir el conteo con ese patrón. \newline \emph{Artefacto:} hoja \emph{Conteos\_base} del Anexo A, con la procedencia de cada conteo.\\ \hline
\end{longtable}

\subsection*{Ética}

\begin{longtable}{|C{0.7cm}|L{14.6cm}|}
\hline
\rowcolor{grisclaro}\textbf{21} & \textbf{¿Qué datos personales trata el sistema y cómo se protegen?}\\ \hline
& Datos identificativos y de contacto del personal, vinculación laboral, registros de avance atribuidos a persona, geolocalización de las capturas e identidad de la cuadrilla. Las medidas son cuatro: cifrado en reposo de la geolocalización con acceso registrado en bitácora (\id{RNF-14}); credenciales almacenadas mediante función de derivación de clave con sal (\id{RNF-13}); disociación de los campos personales antes de incorporar cualquier dato a un conjunto de entrenamiento; y los tres derechos del titular implementados como requisitos funcionales y no como declaraciones: \id{RF-40} acceso, \id{RF-41} rectificación con bitácora y \id{RF-42} supresión con disociación del histórico productivo. Los tres no existían en la Entrega 2A: el ERS enunciaba las obligaciones sin implementarlas, y la inspección lo registró como \id{DEF-12}. \newline \emph{Artefacto:} \S3.4 del ERS; \S\ref{sec:ia}, tabla de base legal.\\ \hline
\rowcolor{grisclaro}\textbf{22} & \textbf{¿Qué partes del informe fueron asistidas por IA y cómo las validaron?}\\ \hline
& Cuatro tareas, declaradas en el Anexo E: la extracción programática de los conteos base, la redacción de las especificaciones de \id{CU-11} a \id{CU-18} y de los criterios \id{CB-01} a \id{CB-37}, la construcción de los diagramas de flujo de datos y de la matriz extremo a extremo, y la redacción de la \S9 sobre componentes de IA. La validación fue distinta en cada caso: los conteos se reprodujeron por muestreo con los patrones declarados; los criterios se contrastaron contra el umbral y la unidad que la ficha ya declaraba, comprobando que ninguno introduce valores nuevos; la correspondencia entre procesos, bloques y requisitos se verificó contra la \S3.2 del ERS; y la clasificación de riesgo se contrastó contra el articulado del Reglamento. No fueron asistidas la retrospectiva, el aporte individual ni las decisiones del CCB. \newline \emph{Artefacto:} Anexo E, declaración de uso de IA.\\ \hline
\end{longtable}

\section{Aporte individual verificable}

\begin{table}[H]
\centering\footnotesize
\caption{Aporte individual por integrante}
\label{tab:aporte}
\begin{tabular}{|L{3.1cm}|L{5.2cm}|C{1.5cm}|L{4.8cm}|}
\hline
\rowcolor{grisclaro}\textbf{Integrante} & \textbf{Aporte principal en la Unidad V} & \textbf{Commits} & \textbf{Artefactos}\\ \hline
Villafuerte Rosero Allan Noe & Consolidación del ERS v1.1 y redacción de la carátula, la \S1 y la \S2 del informe; instrucción de las tres RFC ante el CCB; publicación de la etiqueta de línea base & \pendiente{n.º} & Fuentes \LaTeX{} del ERS y su bibliografía; \S1 y \S2 del informe; acta del CCB y \id{RFC-01} a \id{RFC-03}\\ \hline
Huilcapi León Denisses Fabiola & Estructura editorial del informe: resumen bilingüe, \S3, \S6 y \S10; verificación del formato IEEE de las 51 referencias y redacción de la declaración de uso de IA & \pendiente{n.º} & \S3, \S6 y \S10 del informe; anexos y archivo \id{.bib} de referencias\\ \hline
Rizzo Vélez Edson Nagib & Auditoría de las seis métricas sobre conteos programáticos y re-inspección del ERS corregido; redacción de la \S5, la \S8 y la \S9 del informe & \pendiente{n.º} & Anexo A (auditoría) y Anexo B (registro de defectos) en formato \id{.xlsx}; \S5, \S8 y \S9 del informe\\ \hline
Macías Herrera Josthyn Esteban & Modelado UML del dominio y especificación de los componentes de inteligencia artificial: fichas, umbrales de equidad y explicabilidad, y clasificación de riesgo & \pendiente{n.º} & Diagramas de casos de uso, clases y estados del ERS; \S4 y \S7 del informe\\ \hline
Arboleda Yanza Francisco Javier & Construcción de los diagramas de flujo de datos y de las figuras del informe en TikZ y pgfplots; preparación y reparto por bloques de la presentación de defensa & \pendiente{n.º} & Figuras de la \S8 del ERS; diapositivas de defensa y guion de reparto\\ \hline
Alcívar Vélez Anderson Adonis & Cierre de la trazabilidad extremo a extremo y sincronización del \emph{backlog} con el ERS; estructura del repositorio y redacción del \id{README.md} con las instrucciones de compilación & \pendiente{n.º} & Matriz extremo a extremo y exportación del \emph{backlog} en \id{.xlsx} y \id{.csv}; capturas del tablero; \id{README.md}\\ \hline
\end{tabular}
\end{table}

\noindent
El recuento de commits se obtiene mediante \verb|git shortlog -sne| sobre el repositorio del equipo y corresponde al corte realizado previo al cierre documental. Los commits posteriores asociados a la incorporación de anexos o ajustes documentales no se consideran en el recuento declarado en esta tabla.

\section{Retrospectiva}

La retrospectiva Start--Stop--Continue completa, con sus doce elementos y sus responsables, se
presenta en la \S\ref{sec:retro} del cuerpo del informe.

\section{Declaración de uso de inteligencia artificial}

\begin{longtable}{|L{2.4cm}|L{2.2cm}|L{4.4cm}|L{2.4cm}|L{3.2cm}|}
\hline
\rowcolor{grisclaro}\textbf{Herramienta} & \textbf{Versión} & \textbf{Tarea asistida} & \textbf{Fragmento} & \textbf{Verificación}\\ \hline
\endhead
Claude (Anthropic) & Claude Opus 5 & Extracción programática de los conteos base de la auditoría desde las fuentes \LaTeX{} del ERS & \S\ref{sec:metricas}, Anexo A & El equipo reprodujo por muestreo los conteos con los patrones declarados y comprobó que coinciden\\ \hline
Claude (Anthropic) & Claude Opus 5 & Redacción de las especificaciones de \id{CU-11} a \id{CU-18} y de los criterios \id{CB-01} a \id{CB-37} & \S4 y Apéndice F del ERS & Cada criterio se contrastó contra el umbral y la unidad ya declarados en la ficha del requisito correspondiente; ninguno introduce valores nuevos\\ \hline
Claude (Anthropic) & Claude Opus 5 & Construcción de los diagramas de flujo de datos en TikZ y de la matriz extremo a extremo & \S8 del ERS, \S\ref{sec:gestion} & El equipo verificó la correspondencia entre procesos, bloques funcionales y requisitos contra la \S3.2 del ERS\\ \hline
Claude (Anthropic) & Claude Opus 5 & Redacción de la \S9 del ERS sobre componentes de inteligencia artificial & \S\ref{sec:ia} & Los umbrales se contrastaron con los valores ya especificados en \id{RNF-01} a \id{RNF-05}; la clasificación de riesgo se verificó contra el articulado del Reglamento (UE) 2024/1689\\ \hline
Ninguna & --- & Retrospectiva, aporte individual y decisiones del Change Control Board & \S\ref{sec:retro}, Anexo C & Proceden de artefactos propios del equipo y se redactaron sin asistencia\\ \hline
\caption{Declaración de uso de asistentes de inteligencia artificial}
\label{tab:decl-ia}
\end{longtable}

\vspace{0.4cm}
\noindent\small
El equipo declara que la asistencia se empleó como instrumento de extracción, tabulación y redacción,
y que en ningún caso se incorporó al informe un identificador, una cifra o una referencia que no haya
sido verificada contra el ERS del SIMPA o contra la fuente citada.\normalsize

\vspace{1.2cm}
\noindent\begin{tabular}{p{4.8cm} p{4.8cm} p{4.8cm}}
\hrulefill & \hrulefill & \hrulefill\\
\footnotesize Villafuerte Rosero Allan Noe & \footnotesize Huilcapi León Denisses Fabiola & \footnotesize Rizzo Vélez Edson Nagib\\[1.4cm]
\hrulefill & \hrulefill & \hrulefill\\
\footnotesize Macías Herrera Josthyn Esteban & \footnotesize Arboleda Yanza Francisco Javier & \footnotesize Alcívar Vélez Anderson Adonis\\
\end{tabular}


\end{document}
